\chapter{Local structures}
% archon:covers Proetale/Algebra/WContractible.lean Proetale/Algebra/IdentifiesLocalRings.lean

This chapter serves for the first half of the commutative algebra preparation of the remaining part of the paper. The key result is
\begin{enumerate}
  \item \Cref{thm:ind-etale-w-contractible-cover-exists} expressing that For any ring $A$, there 
  exists an ind-étale faithfully flat $A$-algebra $D$ with $D$ w-contractible. That is to say, 
  every affine scheme has a w-contractible cover in the "pro-(etale site)".
\end{enumerate}

% Question: Can we avoid using weakly etale algebra, only use ind-etale algebra? This would need a modification of theorem 5.4.2, by directly mimicking the proof of topological invariance of étale sites for weakly étale.
% If so, we can remove theorem \Cref{thm:weakly-etale-ind-etale} and every dependency of it.
\section{Preliminaries}

\begin{definition}[Extremally disconnected space]
  \label{def:extremally-disconnected}
  \lean{ExtremallyDisconnected}
  \mathlibok

  A topological space \(X\) is \emph{extremally disconnected} if the closure of every open subset is open.
\end{definition}

The reason that we are interested in extremally disconnected space is the following theorem.

\begin{theorem}
  \label{thm:extremally-disconnected-projective}
  \uses{def:extremally-disconnected}
  \lean{CompHaus.toStonean,Stonean.instProjectiveCompHausCompHaus}
  \mathlibok
  Let \(X\) be a quasi-compact Hausdorff topological space. Then \(X\) is extremally disconnected if and only if it is a projective object in the category of quasi-compact Hausdorff spaces, i.e., for every continuous surjection \(f : Y \to Z\) of quasi-compact Hausdorff spaces and every continuous map \(g : X \to Z\), there exists a continuous map \(h : X \to Y\) such that \(f \circ h = g\).
\end{theorem}

\begin{proof}
  \mathlibok
  Omitted.
\end{proof}

\begin{definition} [Stone-Čech compactification]
  \label{def:stone-cech-compactification}
  \lean{StoneCech}
  \mathlibok
  Let \(X\) be a topological space. The \emph{Stone-Čech compactification} of \(X\) is the
  profinite space \(\beta(X)\) such that
  \begin{enumerate}
    \item \(X\) is dense in \(\beta(X)\);
    \item every continuous map \(f: X \to Y\) to a quasi-compact Hausdorff space \(Y\)
        extends uniquely to a continuous map \(\beta(f): \beta(X) \to Y\).
  \end{enumerate}
  We denote the Stone-Čech compactification of \(X\) by \(\beta(X)\).
\end{definition}

\begin{theorem}
    Let \(X\) be a topological space. Then the Stone-Čech compactification \(\beta(X)\) is extremally disconnected.
    \label{thm:stone-cech-extremally-disconnected}
    \uses{def:extremally-disconnected, def:stone-cech-compactification}
    \lean{StoneCech.projective}
    \mathlibok
\end{theorem}

\begin{proof}
  \mathlibok
  Omitted.
\end{proof}

\begin{proposition}[{\cite[\href{https://stacks.math.columbia.edu/tag/090D}{Tag 090D}]{stacks-project}}]
  Let $X$ be a quasi-compact Hausdorff space. There exists a continuous surjection $X' \to X$ with $X'$ quasi-compact, Hausdorff, and extremally disconnected.
  \label{thm:extremally-disconnected-cover}
  \uses{def:extremally-disconnected, def:stone-cech-compactification}
  \lean{CompHaus.presentation,CompHaus.presentation.π,CompHaus.presentation.epi_π,CompHaus.lift}
  \mathlibok
\end{proposition}

\begin{proof}
  \mathlibok
  \uses{thm:stone-cech-extremally-disconnected}
  Let $Y=X$ but endowed with the discrete topology. Let $X'=\beta (Y)$, which is extremally disconnected by \Cref{thm:stone-cech-extremally-disconnected}. The continuous map $Y \to X$ factors as $Y \to X' \to X$.
\end{proof}

We need the following lemmas about connected components. {\color{red} jiedong: This is WRONG. 
In general, $\pi_0(X) \times \pi_0(Y)$ is continuous and bijective, but not necessarily homeomorphic.}

\begin{lemma}
  \label{thm:pi0-prod}
  \lean{connectedComponent.prod,ConnectedComponents.bijective_connectedComponentsLift_prod,
    ConnectedComponents.isHomeomorph_connectedComponentsLift_prod,
    ConnectedComponents.prodMap}
  Let \(X, Y\) be topological spaces. Then we have 
  \(\pi_0 (X \times Y) \simeq \pi_0(X) \times \pi_0(Y)\) as topological spaces.
\end{lemma}

\begin{proof}
  Let $x, y$ be points of \(X, Y\), respectively. To show the above map is bijective, 
  it suffices to show that the connected component of \((x,y)\) in \(X \times Y\) is the product 
  of the connected component of \(x\) in \(X\) and  the connected component of \(y\) in \(Y\).
  We know that the product of connected subsets is still connected (\verb`IsPreconnected.prod`). 
  The connected component of \((x, y)\) cannot be larger because its projections to \(X\) and \(Y\) 
  are also connected. 

  To show that the bijective map is homeomorphic, we notice that the topology on both side
  are generated by the basis of form \(\text{image } U \times \text{image } V\).
\end{proof}

\begin{lemma}
  \label{thm:pi0-isom}
  \lean{ConnectedComponents.mkHomeomorph}
  \leanok
  Let \(X\) be a totally disconnected topological space. 
  Then \(X \simeq \pi_0 (X)\) as topological spaces.
\end{lemma}

\begin{proof}
  Use the universal property of connected components.
\end{proof}

\begin{lemma}
  Let $X = \lim_i X_i$ be a limit of topological spaces over some index category \(I\) and let
  $s \subseteq X$ be a subset of \(X\). Denote by $s_i$ the image of $s$ in $X_i$. Then
  $\overline{s} = \bigcap_i \pi_i^{-1} (\overline{s_i})$, where $\pi_i : X \to X_i$ is the projection map.
  \label{lemma:closure-limit-intersection}
  \lean{TopCat.closure_eq_iInter_preimage_closure_image}
  \leanok
\end{lemma}

\begin{proof}
  \leanok
  Omitted.
\end{proof}

\begin{lemma}
  \label{thm:closure-limit-compatible}
  Let $X = \lim_i X_i$ be a limit of topological spaces over some index category \(I\) and let
  $s \subseteq X$ be a subset of \(X\). Denote by $s_i$ the image of $s$ in $X_i$. Then for any morphism \(f : i \to j\) in \(I\) with transition map \(\varphi_{f} : X_i \to X_j\), we have \(\varphi_{f}(\overline{s_i}) \subseteq \overline{s_j}\).
  \lean{image_closure_image_subset_closure_image}
  \leanok
\end{lemma}

\begin{proof}
  \leanok
  Omitted.
\end{proof}


\begin{lemma}
    Let $f\colon X \to Y$ be an embedding of topological spaces. Then $f$ preserves and reflects
    specializations, i.e. for $a, b \in X$, $a$ specializes to $b$ if and only if $f(a)$ specializes to $f(b)$.
    \label{lemma:embedding-specializations}
    \lean{Topology.IsEmbedding.specializes_iff}
    \leanok
\end{lemma}

\begin{proof}
    \leanok
    The only if direction holds for any continuous map. Conversely, if $f(a)$ specializes to $f(b)$, then
    $b \in f^{-1}(\overline{\{f(a)\}}) = \overline{\{a\} }$.
\end{proof}

\begin{definition}
    Let $X$ be a topological space and $Z \subseteq X$ a subset. The set
    \[
    Z^g = \{ x \in X \mid \exists z \in Z, x \rightsquigarrow z \}
    \] is called the \emph{generalization of $Z$ in $X$}.
    \label{def:generalization}
    \lean{generalizationHull}
    \leanok
\end{definition}

\begin{lemma}
    Let $Z \subseteq X$ be a subset of a topological space. Then $Z^g$ is closed under generalization.
    \label{lemma:generalization-stable-generalization}
    \lean{stableUnderGeneralization_generalizationHull}
    \leanok
\end{lemma}

\begin{proof}
    \leanok
    This is immediate from the definition.
\end{proof}

The following pure topological lemmas will be used in \Cref{thm:colim-open-closed-localizations-surj-ind-zariski,thm:cartesian-profinite-w-local}.

\begin{definition}[{\cite[\href{https://stacks.math.columbia.edu/tag/08ZL}{Tag 08ZL}]{stacks-project}}]
  \label{def:pi0-to-totally-disconnected}
  \lean{Continuous.connectedComponentsLift,Continuous.connectedComponentsLift_continuous,Continuous.connectedComponentsLift_apply_coe,Continuous.connectedComponentsLift_comp_coe}
  \mathlibok
  Let \(X\) be a topological space. Let \(\pi_0(X)\) be the set of connected components of \(X\). Let \(X \to \pi_0(X)\) be the map which sends \(x \in X\) to the connected component of \(X\) passing through \(x\). Endow \(\pi_0(X)\) with the quotient topology. Then %\(\pi_0(X)\) is a totally disconnected space and 
  any continuous map \(X \to Y\) from \(X\) to a totally disconnected space \(Y\) factors through \(\pi_0(X)\).
\end{definition}

\begin{lemma}
  \label{thm:pi0-to-totally-disconnected-injective}
  \uses{def:pi0-to-totally-disconnected}
  \lean{Continuous.connectedComponentsLift_injective}
  \leanok
  Let \(X\) be a topological space. Let \(\pi_0(X)\) be the topological space of connected components of \(X\). Let \(Y\) be a totally disconnected space. If every fiber of \(X \to Y\) is connected, then the map \(\pi_0(X) \to Y\) is injective.
\end{lemma}

\begin{proof}
  Omitted.
\end{proof}

\begin{lemma}[{\cite[\href{https://stacks.math.columbia.edu/tag/005F}{Tag 005F}]{stacks-project}}]
  \label{thm:connected-component-intersection-closed-open}
  \lean{sInter_isClopen_and_mem_eq_connectedComponent}
  \leanok
  Let $X$ be a topological space. Assume $X$ is quasi-compact, quasi-separated
  and prespectral. Then for any $x \in X$ the connected component of $X$ containing $x$ is
  the intersection of all open and closed subsets of $X$ containing $x$.
\end{lemma}

\begin{proof}
    \leanok
  Let $T$ be the connected component containing $x$. Let $S = \bigcap_{\alpha \in A} Z_\alpha$ be the intersection of all open and closed subsets $Z_\alpha$ of $X$ containing $x$. Note that $S$ is closed in $X$. Note that any finite intersection of $Z_\alpha$'s is a $Z_\alpha$. Because $T$ is connected and $x \in T$ we have $T \subset Z_\alpha$ for every \(Z_\alpha\), thus \(T \subset S\). It suffices to show that $S$ is connected. If not, then there exists a disjoint union decomposition $S = B \amalg C$ with $B$ and $C$ open and closed in $S$. In particular, $B$ and $C$ are closed in $X$, and so quasi-compact by assumption (1). By assumption (2) there exist quasi-compact opens $U, V \subset X$ with $B = S \cap U$ and $C = S \cap V$. Then $U \cap V \cap S = \emptyset$. Hence $\bigcap_\alpha U \cap V \cap Z_\alpha = \emptyset$. By assumption (3) the intersection $U \cap V$ is quasi-compact. Hence for some $\alpha' \in A$ we have $U \cap V \cap Z_{\alpha'} = \emptyset$. Since $X - (U \cup V)$ is disjoint from $S$ and closed in $X$ hence quasi-compact, we can use the same argument to see that $Z_{\alpha''} \subset U \cup V$ for some $\alpha'' \in A$. Then $Z_\alpha = Z_{\alpha'} \cap Z_{\alpha''}$ is contained in $U \cup V$ and disjoint from $U \cap V$. Hence $Z_\alpha = U \cap Z_\alpha \amalg V \cap Z_\alpha$ is a decomposition into two open pieces, hence $U \cap Z_\alpha$ and $V \cap Z_\alpha$ are open and closed in $X$. Thus, if $x \in B$ say, then we see that $S \subset U \cap Z_\alpha$ and we conclude that $C = \emptyset$.
\end{proof}

\begin{lemma}[{\cite[\href{https://stacks.math.columbia.edu/tag/04PL}{Tag 04PL}]{stacks-project}}]
  \label{thm:intersection-closed-open-iff}
  \lean{isClosed_and_iUnion_connectedComponent_eq_iff}
  \leanok
  Let $X$ be a topological space. Assume $X$ is quasi-compact, quasi-separated
  and prespectral. Then for a subset $T \subset X$ the following are equivalent:
  \begin{enumerate}
    \item $T$ is an intersection of open and closed subsets of $X$, and
    \item $T$ is closed in $X$ and is a union of connected components of $X$.
  \end{enumerate}
\end{lemma}

\begin{proof}
  \leanok
  \uses{thm:connected-component-intersection-closed-open}
  It is clear that (a) implies (b). Assume (b). Let $x \in X$, $x \notin T$. Let $x \in C \subset X$ be the connected component of $X$ containing $x$. By \Cref{thm:connected-component-intersection-closed-open} we see that $C = \bigcap V_\alpha$ is the intersection of all open and closed subsets $V_\alpha$ of $X$ which contain $C$. In particular, any pairwise intersection $V_\alpha \cap V_\beta$ occurs as a $V_\alpha$ i.e. is still open and closed. As $T$ is a union of connected components of $X$ we see that $C \cap T = \emptyset$. Hence $T \cap \bigcap V_\alpha = \emptyset$. Since $T$ is quasi-compact as a closed subset of a quasi-compact space, we deduce that $T \cap V_\alpha = \emptyset$ for some $\alpha$. For this $\alpha$ we see that $U_\alpha = X - V_\alpha$ is an open and closed subset of $X$ which contains $T$ and not $x$. The lemma follows.
\end{proof}

\begin{lemma}[{\cite[\href{https://stacks.math.columbia.edu/tag/0906}{Tag 0906}]{stacks-project}}]
  \label{thm:pi0-profinite}
  \lean{compactSpace_connectedComponent,t2Space_connectedComponent}
  Let $X$ be a spectral space. Then $\pi_0(X)$ is a profinite space.
\end{lemma}

\begin{proof}
  \uses{thm:connected-component-intersection-closed-open}
  We will show that \(\pi_0(X)\) is Hausdorff, quasi-compact, and totally disconnected.
  By \Cref{thm:connected-component-intersection-closed-open}, every connected component of \(X\) is the intersection of the open and closed subsets containing it. Since \(\pi_0(X)\) is the image of a quasi-compact space, it is also quasi-compact. It is totally disconnected by construction. Let \(C,D \subset X\) be distinct connected components of \(X\). Write \(C = \bigcap_{\alpha} U_\alpha\) as the intersection of the open and closed subsets of \(X\) containing \(C\). Any finite intersection of \(U_\alpha\)'s is another \(U_\alpha\). Since \(\bigcap_\alpha U_\alpha \cap D = \emptyset\) we conclude that \(U_\alpha \cap D = \emptyset\) for some \(\alpha\) (connected components are closed, thus quasi-compact). Since \(U_\alpha\) is open and closed, it is the union of the connected components it contains, i.e., \(U_\alpha\) is the inverse image of some open and closed subset \(V_\alpha \subset \pi_0(X)\). This proves that the points corresponding to \(C\) and \(D\) are contained in disjoint open subsets, i.e., \(\pi_0(X)\) is Hausdorff.
\end{proof}

\begin{lemma}
  \label{thm:closed-subspace-spectral}
  \lean{Topology.IsClosedEmbedding.isOpen_and_isCompact_and_preimage_eq,
    Topology.IsClosedEmbedding.spectralSpace,SpectralSpace.of_isClosed}
  Let $X$ be a spectral space. Let $Y \subset X$ be a closed subspace. Then $Y$ is spectral.
\end{lemma}

\begin{proof}
  The only remaining part in Lean is quasi-separatedness.
  Suppose that \(V\) is open compact in \(Y\), we show that there exists open compact \(U\) in 
  \(X\) such that \(V = U \cap Y\). We can find \(U'\) in \(X\) such that \(V = U' \cap Y\). Write
  \(U' = \bigcup_{i \in I} U_i\) as a union of compact opens, then a finite subset \(i \in I_0\) 
  of \(U_i\) already covers \(V\). Take \(U := \bigcup_{i \in I\_0} U_i\) as desired.
  
  For any compact opens \(V_1,\, V_2 \subset Y\), 
  there exists compact open \(U_1, U_2 \subset X\) s.t. \(V_i = U_i \cap Y_i\). Then 
  \(V_1 \cap V_2 = (U_1 \cap U_2) \cap Y \), is still open compact in \(Y\). Here \(U_1 \cap U_2\)
  is quasi-compact since \(X\) is quasi-separated.
\end{proof}

\begin{lemma}
  \label{thm:closure-product}
  \lean{closure_prod_eq,Set.singleton_prod_singleton}
  \mathlibok
  Let $X$ and $Y$ be topological spaces. Let $(x, y) \in X \times Y$ be a point. Then 
  \[\overline{\{(x, y)\}} = \overline{\{x\}} \times \overline{\{y\}}.\]
\end{lemma}

\begin{proof}
  \mathlibok
  It is clear that \(\overline{\{(x, y)\}} \subseteq \overline{\{x\}} \times \overline{\{y\}}\).

  If \(z \notin \overline{\{(x, y)\}}\), there exists opens \(U \subset X\), \(V \subset Y\) such that \(z \notin U \times V\), thus \((x, y) \notin U \times V\). Either \(x \notin U\) or \(y \notin V\), suppose \(x \notin U\). Then \(U \cap \overline{\{x\}} = \emptyset\), thus \(U \times V \cap \overline{\{x\}} \times \overline{\{y\}}  = \emptyset\). This shows that \(\overline{\{x\}} \times \overline{\{y\}} \subseteq \overline{\{(x, y)\}}\).
\end{proof}

\begin{lemma}[{\cite[\href{https://stacks.math.columbia.edu/tag/0907}{Tag 0907}]{stacks-project}}]
  \label{thm:spectral-product}
  \lean{QuasiSeparatedSpace.prod,QuasiSober.prod,PrespectralSpace.prod,SpectralSpace.prod}
  Let $X$ and $Y$ be spectral spaces. Then the product $X \times Y$ is spectral.
\end{lemma}

\begin{proof}
  \uses{thm:closure-product}
  Let $X$, $Y$ be spectral spaces. Denote $p : X \times Y \to X$ and $q : X \times Y \to Y$ the projections. We first show that \(X \times Y\) is sober. Let $Z \subset X \times Y$ be a closed irreducible subset. Then $p(Z) \subset X$ is irreducible and $q(Z) \subset Y$ is irreducible. Let $x \in X$ be the generic point of the closure of $p(Z)$ and let $y \in Y$ be the generic point of the closure of $q(Z)$. If $(x, y) \notin Z$, then there exist opens $x \in U \subset X$, $y \in V \subset Y$ such that $Z \cap (U \times V) = \emptyset$. Hence $Z$ is contained in $(X - U) \times Y \cup X \times (Y - V)$. Since $Z$ is irreducible, we see that either $Z \subset (X - U) \times Y$ or $Z \subset X \times (Y - V)$. In the first case $p(Z) \subset (X - U)$ and in the second case $q(Z) \subset (Y - V)$. Both cases are absurd as $x$ is in the closure of $p(Z)$ and $y$ is in the closure of $q(Z)$. Thus we conclude that $(x, y) \in Z$, which means that $(x, y)$ is the generic point for $Z$ (\Cref{thm:closure-product}).

  A basis of the topology of $X \times Y$ consists of opens of the form $U \times V$ with $U \subset X$ and $V \subset Y$ quasi-compact open (here we use that $X$ and $Y$ are spectral). Then $U \times V$ is quasi-compact as the product of quasi-compact spaces is quasi-compact. Moreover, any quasi-compact open of $X \times Y$ is a finite union of such quasi-compact rectangles $U \times V$. It follows that the intersection of two such is again quasi-compact (since $X$ and $Y$ are spectral). This concludes the proof.
\end{proof}

\begin{lemma}[{\stacksproject{08ZE}}]
  \label{thm:Hausdorff-diagonal-closed}
  \lean{t2Space_iff_diagonal_closed}
  \leanok
  Let \(X\) be a topological space. Then \(X\) is Hausdorff if and only if the diagonal map \(\Delta : X \to X \times X\) is a closed embedding.
\end{lemma}

\begin{proof}
  \leanok
  Omitted.
\end{proof}

\begin{lemma}[{\stacksproject{08ZH}}]
  \label{thm:pullback-closed-embedding-prod}
  \lean{TopCat.isClosedEmbedding_pullback_to_prod}
  Let \(X\), \(Y\), and \(S\) be topological spaces. Assume that \(S\) is Hausdorff. 
  Let \(f : X \to S\) and \(g : Y \to S\) be continuous maps. The natural map 
  \(X \times_S Y \to X \to Y\) is a closed embedding.
\end{lemma}

\begin{proof}
  \uses{thm:Hausdorff-diagonal-closed}
  \(X \times_S Y\) is the inverse image of the image \(\Delta(S)\) of the diagonal map 
  \(\Delta : S \to S \times S\). It follows from \Cref{thm:Hausdorff-diagonal-closed} that 
  \(\Delta(S)\) is a closed subset.
\end{proof}

\begin{lemma}
  \label{thm:homeo-preserves-spectral}
  \lean{Homeomorph.spectralSpace,Homeomorph.t0Space,Homeomorph.t2Space,Homeomorph.quasiSober,Homeomorph.quasiSeparatedSpace,Homeomorph.compactSpace,Homeomorph.prespectralSpace}
  Let \(X\) and \(Y\) be topological spaces. If \(X\) is spectral (resp. T0, Hausdorff, Sober, quasi-separated, quasi-compact, prespectral) and there is a homeomorphism \(X \cong Y\), then \(Y\) is spectral (resp. T0, Hausdorff, Sober, quasi-separated, quasi-compact, prespectral).
\end{lemma}

\begin{proof}
  Omitted.
\end{proof}

\begin{lemma}
  \label{thm:pullback-spectral} 
  \lean{T0Space.pullback,PrespectralSpace.pullback,SpectralSpace.pullback,
    QuasiSober.pullback,CompactSpace.pullback,CategoryTheory.IsPullback.t0Space,
    CategoryTheory.IsPullback.quasiSober,CategoryTheory.IsPullback.compactSpace,
    CategoryTheory.IsPullback.prespectralSpace,CategoryTheory.IsPullback.spectralSpace}
  Let \(X\) and \(Y\) be spectral spaces. Let \(S\) be a Hausdorff topological space. Let \(f : X \to S\) and \(g : Y \to S\) be continuous maps. Then the fiber product \(X \times_S Y\) is spectral.
\end{lemma}

\begin{proof}
  \uses{thm:spectral-product,thm:closed-subspace-spectral,thm:pullback-closed-embedding-prod,
    thm:homeo-preserves-spectral}
  Note that \(X \times_S Y\) is a closed subspace of $X \times Y$ by 
  \Cref{thm:pullback-closed-embedding-prod}. (For subspace, see 
  \verb`TopCat.pullbackIsoProdSubtype`.) Since $X \times Y$ is spectral 
  (\Cref{thm:spectral-product}) it follows that \(X \times_S Y\) is spectral 
  (\Cref{thm:closed-subspace-spectral}). Note that we also used \Cref{thm:homeo-preserves-spectral}
  to translate between categorical constructions and non-categorical constructions.
\end{proof}

\begin{lemma}
  \label{thm:profinite-spectral}
  \lean{Prespectral.of_profinite,Spectral.of_profinite}
  Let $X$ be a profinite space. Then $X$ is spectral.
\end{lemma}

\begin{proof}
  Since R1 implies Sober, Hausdorff implies quasi-separated, it only remains to show that \(X\) is prespectral. Write \(X\) as an inverse limit of finite discrete topological spaces. The preimage of singletons in every finite discrete space are open and closed subsets and forms a topological basis for the topology on \(X\). Thus \(X\) is prespectral.
\end{proof}

\begin{lemma}[{\cite[\href{https://stacks.math.columbia.edu/tag/096C}{Tag 096C}, first part]{stacks-project}}]
\label{thm:cartesian-profinite}
\uses{def:pi0-to-totally-disconnected}
\lean{ConnectedComponents.spectralSpace_of_isPullback,ConnectedComponents.lift_bijective_of_isPullback,ConnectedComponents.isHomeomorph_lift_of_isPullback}
Let $X$ be a spectral space. Let
\[
\begin{tikzcd}
Y \arrow[r] \arrow[d] & T \arrow[d] \\
X \arrow[r] & \pi_0(X)
\end{tikzcd}
\]
be a cartesian diagram in the category of topological spaces with $T$ profinite. Then $Y$ is spectral and the canonical map $\pi_0(Y) \to T$ defined in \Cref{def:pi0-to-totally-disconnected} is an isomorphism.
\end{lemma}

\begin{proof}
  \uses{thm:profinite-spectral,thm:pullback-spectral,thm:pi0-profinite,def:pi0-to-totally-disconnected,thm:pi0-to-totally-disconnected-injective}
  By \Cref{thm:profinite-spectral}, \(T\) is spectral. By \Cref{thm:pi0-profinite} \(\pi_0(X)\) is Hausdorff.
  By \Cref{thm:pullback-spectral}, \(Y\) is spectral.
  Let $Y \to \pi_0(Y) \to T$ be the canonical factorization in \Cref{def:pi0-to-totally-disconnected}. It is clear that $\pi_0(Y) \to T$ is surjective. The fibres of $Y \to T$ are homeomorphic to the fibres of $X \to \pi_0(X)$. Hence these fibres are connected. It follows that $\pi_0(Y) \to T$ is injective by \Cref{thm:pi0-to-totally-disconnected-injective}. We conclude that $\pi_0(Y) \to T$ is a homeomorphism since it is a bijective map from a quasi-compact space to a Hausdorff space by \verb`isHomeomorph_iff_continuous_bijective`.
\end{proof}

The following pure algebraic lemma will be used in \Cref{thm:isom-of-identifies-local-rings-w-local-isom-pi0}.

\begin{lemma}[{\cite[\href{https://stacks.math.columbia.edu/tag/00EA}{Tag 00EA}, (1) + (2)(a)]{stacks-project}}]
  \label{thm:localization-isom-of-going-down}
  \lean{Algebra.HasGoingDown.isLocalization_atPrime_of_subsingleton}
  \leanok
  Let $A \to B$ be a ring map which has going down. Let $\p \subset A$ be a prime ideal and let 
  $\q \subset B$ be a prime ideal lying over $\p$. Assume that and there is at most one prime
  of $B$ above every prime of $A$. Then the natural map $B_{\p B} \to B_{\q}$ is an isomorphism.
\end{lemma}

\begin{proof}
  \leanok
  It suffices to show that all the elements of $B - \q$ are all invertible in $B_{\p B}$. Suppose the contrary, let $x \in B - \q$ be an element not invertible in $B_{\p B}$. We can find a prime ideal $\q' \subseteq \p B$ of $B$ containing $x$. By going down there exists a prime $\q'' \subseteq \q$ lying over $\p' = A \cap \q'$. By the uniqueness of primes lying over $\p'$ we see that $\q' = \q''$, which contradicts the fact that $x \notin \q$.
\end{proof}

\begin{definition}[Local isomorphisms, {\cite[\href{https://stacks.math.columbia.edu/tag/096E}{Tag 096E} (1)]{stacks-project}}]
  \label{def:local-isomorphism}
  \lean{Algebra.IsLocalIso, RingHom.IsLocalIso}
  We say $A \to B$ is a \emph{local isomorphism} if for every prime $\mathfrak{q} \subset B$ there exists a $g \in B$, $g \notin \mathfrak{q}$ such that $A \to B_g$ induces an open immersion $\Spec (B_g) \to \Spec (A)$.
\end{definition}

\begin{definition}[Identify local rings, {\cite[\href{https://stacks.math.columbia.edu/tag/096E}{Tag 096E} (2)]{stacks-project}}]
  A ring map $A \to B$ \emph{identifies local rings} if for every prime $\mathfrak{q} \subset B$ the canonical map $A_{\varphi^{-1}(\mathfrak{q})} \to B_{\mathfrak{q}}$ is an isomorphism.
  \label{def:identify-local-rings}
  \lean{RingHom.BijectiveOnStalks}
  \leanok
\end{definition}

\begin{lemma}
  \label{thm:finite-product-identifies-local-rings}
  \uses{def:identify-local-rings}
  \lean{RingHom.BijectiveOnStalks.prod}
  \leanok
  The map \(A \to \prod_{i = 1}^{n} A_i \) identifies local rings, if \(A \to A_i\) identifies
  local rings for all \(i\).
\end{lemma}

\begin{proof}

\leanok
  Omitted.
\end{proof}

\begin{lemma}
  \label{thm:identifies-local-rings-composition}
  Let \(A \to B\) and \(B \to C\) be ring maps that identify local rings. Then the composition \(A \to C\) also identifies local rings.
  \uses{def:identify-local-rings}
  \lean{RingHom.BijectiveOnStalks.comp}
  \leanok
\end{lemma}

\begin{proof}
  \leanok
  Omitted.
\end{proof}

\begin{lemma}[\stacksproject{096D}]
    Let $A$ be a ring and $B \to C$ be an $A$-algebra map. Suppose both
    $A \to B$ and $A \to C$ are bijective on stalks. Then $B \to C$ is bijective on stalks.
    \label{lemma:bijective-on-stalks-of-comp}
\end{lemma}

\begin{proof}
    Holds because the same holds for bijective maps of sets.
\end{proof}

\begin{definition}[{\cite[\href{https://stacks.math.columbia.edu/tag/096L}{Tag 096L}]{stacks-project}}]
  \label{def:identifies-local-ring-to-top}
  \uses{def:identify-local-rings,thm:identifies-local-rings-composition}
  Let $A$ be a ring and $X = \Spec(A)$. Let $\operatorname{ILR}_A$ denote the category of $A$-algebras $B$ for which $A \to B$ identifies local rings, and let $\operatorname{Top}_X$ denote the category of topological spaces over $X$.
  
  Define the functor
  \[
  F \colon \operatorname{ILR}_A \longrightarrow \operatorname{Top}_X
  \]
  by sending $B$ to $\Spec(B)$.
\end{definition}

\begin{definition}\leanok
[Lean carrier for ``Hom over $\Spec A$'' --- Archon-original]
  \label{def:identifies-local-ring-homover}
  \uses{def:identify-local-rings}
  \lean{Algebra.BijectiveOnStalks.HomOver}
  Let $A$ be a ring, $X = \Spec(A)$, and let $B, C$ be $A$-algebras with
  structure maps $\alpha_B : A \to B$ and $\alpha_C : A \to C$. Define
  $\mathrm{HomOver}(A, B, C)$ to be the set of pairs $(\varphi, h)$
  where $\varphi : \Spec(C) \to \Spec(B)$ is a continuous map and $h$
  is a witness of the compatibility condition
  $\Spec(\alpha_B) \circ \varphi = \Spec(\alpha_C)$, i.e. for every
  $\mathfrak p \in \Spec(C)$ one has
  $\alpha_B^{-1}(\varphi(\mathfrak p)) = \alpha_C^{-1}(\mathfrak p)$
  as elements of $\Spec(A)$.

  This concretely materialises $\Hom_{\mathrm{Top}_X}(\Spec C, \Spec B)$ as a
  set --- the Lean carrier \texttt{Algebra.BijectiveOnStalks.HomOver}
  is the iter-132 design choice (a concrete \texttt{structure} rather
  than a \texttt{CategoryTheory.Over} object) made to facilitate
  \texttt{FunLike} / \texttt{@[ext]} API.
\end{definition}

\begin{lemma}\leanok
{\cite[\href{https://stacks.math.columbia.edu/tag/096L}{Tag 096L}]{stacks-project}}
  \label{thm:identifies-local-ring-to-top-fully-faithful}
  \lean{Algebra.BijectiveOnStalks.algHomEquivContinuousMap}
  \uses{def:identifies-local-ring-to-top}
  Let $A$ be a ring. Set $X = \Spec (A)$. The functor \(F\) constructed in \Cref{def:identifies-local-ring-to-top}
  \[
  B \longmapsto \Spec (B)
  \]
  from the category of $A$-algebras $B$ such that $A \to B$ identifies local rings to the category of topological spaces over $X$ is fully faithful.

  \emph{Lean delivery (iter-132).} The Lean carrier lives in
  \texttt{Proetale/Algebra/IdentifiesLocalRings.lean} as a hom-set
  bijection (rather than a full categorical functor) — when $A \to B$
  and $A \to C$ both identify local rings (Lean class
  \texttt{Algebra.BijectiveOnStalks}, already present in
  \texttt{Proetale/Algebra/StalkIso.lean}), the map
  $\Hom_{A\text{-alg}}(B, C) \to \mathrm{Hom}_{\mathrm{Top}_X}(\Spec C, \Spec B)$
  is an equivalence of sets. The iter-132 file-skeleton lane scaffolds
  the statement; the closing proof in subsequent iters formalises the
  ringed-space identification $\mathcal O_Y = p^{-1} \mathcal O_X$ for
  $A \to B$ identifying local rings, then invokes the canonical
  isomorphism $f^{-1}\mathcal O_Y \cong q^{-1}\mathcal O_X \cong
  \mathcal O_Z$ on the codomain.
\end{lemma}

% SOURCE QUOTE PROOF: "This follows from Lemma
% \ref{lemma-fully-faithful-spaces-over-X} and the fact that if $A \to B$
% identifies local rings, then the pullback of the structure sheaf of
% $\Spec(A)$ via $p : \Spec(B) \to \Spec(A)$ is equal to the structure
% sheaf of $\Spec(B)$." (Stacks Project Tag 096L, proof, read from
% references/stacks-096J.tex L140-145).
\begin{proof}
\leanok
\uses{def:identify-local-rings,def:identifies-local-ring-homover,def:identifies-local-ring-to-top,
      lem:identifies-local-ring-spec-injective,lem:identifies-local-ring-spec-existence}
\textit{Source: Stacks Project, Tag 096L, proof.}
The functor \(F\) factors as the composition
\[
  \operatorname{ILR}_A \;\longrightarrow\;
    \operatorname{LRS}_{/X} \;\longrightarrow\;
    \operatorname{Top}_{/X},
\]
where the first arrow sends \(B\) to \((\Spec B, \mathcal O_{\Spec B}) \to (\Spec A, \mathcal O_{\Spec A})\) and the second forgets the sheaf-level data. The hypothesis that \(A \to B\) identifies local rings is equivalent to the assertion that the pullback of the structure sheaf is the structure sheaf, i.e.\ \(p^{-1}\mathcal O_X = \mathcal O_Y\) for \(p : Y = \Spec B \to X = \Spec A\); under this hypothesis \Cref{lem:identifies-local-ring-spec-spaces-over-X} shows that the second arrow is bijective on hom-sets. The first arrow is bijective on hom-sets by the affine-scheme \(\Spec/\Gamma\) adjunction.

\medskip

\noindent\emph{Lean-side decomposition.} On the formal side the target
is the equivalence
\[
  \mathrm{algHomEquivContinuousMap} \colon
    \Hom_{A\text{-alg}}(B, C) \;\simeq\; \mathrm{HomOver}(A, B, C),
\]
where the right-hand carrier was introduced in
\Cref{def:identifies-local-ring-homover}. Following the iter-133/134
proof split, we package this as the pair
\((\Phi, \Psi)\) of forward and inverse maps and read off the bijection
from the two halves separately: injectivity of \(\Phi\) is handled in
\Cref{lem:identifies-local-ring-spec-injective}, and existence of a
preimage under \(\Phi\) (i.e.\ surjectivity) is handled in
\Cref{lem:identifies-local-ring-spec-existence}. The bundled equivalence
\(\mathrm{algHomEquivContinuousMap}\) is then assembled by
\texttt{Equiv.ofBijective} from the bijection
\(\Hom_{A\text{-alg}}(B, C) \to \mathrm{HomOver}(A, B, C)\) of
\Cref{lem:identifies-local-ring-spec-injective} +
\Cref{lem:identifies-local-ring-spec-existence}.
\end{proof}

% ------------------------------------------------------------------
% iter-135 split: the iter-133/134 prover decomposed
% thm:identifies-local-ring-to-top-fully-faithful into the LRS->Top
% bijection (Stacks 096K, recorded below as a chapter sub-lemma) plus
% injectivity and existence halves at the AlgHom <-> HomOver level. The
% three sub-lemmas below correspond to the Lean declarations
%   Algebra.BijectiveOnStalks.continuousMap_of_algHom_injective
%   Algebra.BijectiveOnStalks.exists_algHom_of_continuousMap
% (the closed/open halves named in
% .archon/task_results/Algebra_IdentifiesLocalRings.lean.md) plus an
% explicit chapter record of Stacks 096K used in the parent's proof above.
% ------------------------------------------------------------------

% SOURCE: Stacks Project, Tag 096K, Lemma 61.3.7
% (read from references/stacks-096J.tex L108-127).
% SOURCE QUOTE: "Let $p : (Y, \mathcal{O}_Y) \to (X, \mathcal{O}_X)$ and
% $q : (Z, \mathcal{O}_Z) \to (X, \mathcal{O}_X)$ be morphisms of locally
% ringed spaces. If $\mathcal{O}_Y = p^{-1}\mathcal{O}_X$, then
% $$\Mor_{\text{LRS}/(X, \mathcal{O}_X)}((Z, \mathcal{O}_Z), (Y, \mathcal{O}_Y))
% \longrightarrow \Mor_{\textit{Top}/X}(Z, Y),\quad (f, f^\sharp) \longmapsto f$$
% is bijective. Here $\text{LRS}/(X, \mathcal{O}_X)$ is the category of
% locally ringed spaces over $X$ and $\textit{Top}/X$ is the category of
% topological spaces over $X$."
\begin{lemma}[Stacks 096K: LRS-over-$X$ forget-to-Top is bijective when $\mathcal O_Y = p^{-1}\mathcal O_X$]
  \label{lem:identifies-local-ring-spec-spaces-over-X}
  Let \((X, \mathcal O_X)\) be a locally ringed space and let
  \(p \colon (Y, \mathcal O_Y) \to (X, \mathcal O_X)\) and
  \(q \colon (Z, \mathcal O_Z) \to (X, \mathcal O_X)\) be morphisms of
  locally ringed spaces. If \(\mathcal O_Y = p^{-1}\mathcal O_X\), then
  the forgetful map
  \[
    \Mor_{\mathrm{LRS}/(X, \mathcal O_X)}\!\left((Z, \mathcal O_Z), (Y, \mathcal O_Y)\right)
      \;\longrightarrow\;
      \Mor_{\mathrm{Top}/X}(Z, Y), \quad (f, f^\sharp) \longmapsto f
  \]
  is a bijection.
  \textit{Source: Stacks Project, Tag 096K.}
\end{lemma}

% SOURCE QUOTE PROOF: "This is immediate from the definitions." (Stacks
% Project Tag 096K, proof, read from references/stacks-096J.tex L125-127).
\begin{proof}
  \textit{Source: Stacks Project, Tag 096K, proof.}
  Given a continuous \(f \colon Z \to Y\) with \(p \circ f = q\),
  inverse-image functoriality and the hypothesis \(\mathcal O_Y =
  p^{-1}\mathcal O_X\) give a canonical chain
  \[
    f^{-1}\mathcal O_Y \;=\; f^{-1} p^{-1} \mathcal O_X
      \;=\; (p \circ f)^{-1} \mathcal O_X
      \;=\; q^{-1} \mathcal O_X \;=\; \mathcal O_Z;
  \]
  one checks the composed map is an LRS-morphism (the stalk components
  are isomorphisms because \(\mathcal O_Y = p^{-1}\mathcal O_X\)
  identifies each stalk \((\mathcal O_Y)_{f(z)}\) with \((\mathcal O_X)_{p(f(z))} = (\mathcal O_X)_{q(z)} = (\mathcal O_Z)_z\)).
  Conversely, an LRS-morphism over \((X, \mathcal O_X)\) determines
  its underlying continuous map and is determined by it because the
  sheaf component is forced by the same inverse-image chain.
\end{proof}

\begin{lemma}\leanok
[Injectivity of $\Spec$ on $\Hom_{A\text{-alg}}$ in the \texttt{BijectiveOnStalks} setting --- Archon split]
  \label{lem:identifies-local-ring-spec-injective}
  \lean{Algebra.BijectiveOnStalks.continuousMap_of_algHom_injective}
  \uses{def:identify-local-rings,def:identifies-local-ring-homover}
  Let \(A\) be a ring and let \(A \to B\), \(A \to C\) be two
  \(A\)-algebra structures both identifying local rings. The forward
  map
  \[
    \Phi \colon \Hom_{A\text{-alg}}(B, C) \;\longrightarrow\;
      \mathrm{HomOver}(A, B, C),
      \qquad f \longmapsto (\Spec(f), \text{compatibility witness}),
  \]
  introduced in \Cref{def:identifies-local-ring-hom-set-bijection}, is
  injective.
\end{lemma}

\begin{proof}
  \uses{def:identify-local-rings,def:identifies-local-ring-homover}
  \leanok
  Suppose \(f_1, f_2 : B \to C\) are two \(A\)-algebra homomorphisms with
  \(\Spec(f_1) = \Spec(f_2)\). For each prime \(\mathfrak p \subset C\),
  write \(\mathfrak q_i = f_i^{-1}(\mathfrak p)\); the equality of
  Spec-maps gives \(\mathfrak q_1 = \mathfrak q_2\), call this common
  prime \(\mathfrak q\). Localising at \(\mathfrak p\) yields two
  ring maps
  \(M_i \colon B_\mathfrak q \to C_\mathfrak p\)
  induced by \(f_i\), and these agree after pre-composition with the
  canonical \(A_{\mathfrak q \cap A} \to B_\mathfrak q\): both equal the
  local ring map \(A_{\mathfrak q \cap A} \to C_\mathfrak p\)
  attached to \(A \to C\) at \(\mathfrak p\) (this uses that each
  \(f_i\) is an \(A\)-algebra map, so it commutes with the structure
  maps from \(A\)). Because \(A \to B\) identifies local rings, the
  canonical map \(A_{\mathfrak q \cap A} \to B_\mathfrak q\) is
  surjective (in fact an isomorphism), so \(M_1 = M_2\). Finally, the
  product of localisation maps
  \(C \to \prod_{\mathfrak p}\, C_\mathfrak p\) is injective, so
  \(f_1 = f_2\).

  \smallskip

  \noindent\emph{Lean delivery.} The injectivity half lives in
  \texttt{Proetale/Algebra/IdentifiesLocalRings.lean} as
  \texttt{continuousMap\_of\_algHom\_injective}. The chain of
  reductions above uses, in order, \texttt{AlgHom.coe\_ringHom\_injective}
  / \texttt{RingHom.ext} to reduce to pointwise equality, then
  \texttt{PrimeSpectrum.toPiLocalization\_injective C} for the final
  localisation-injectivity step, and \texttt{Localization.localRingHom}
  for the per-prime maps \(M_i\). The proof was closed sorry-free in
  iter-134.

  \smallskip

  \noindent\emph{Source.} Per Stacks 096L the parent functor
  \(\operatorname{ILR}_A \to \operatorname{Top}_X\) is fully faithful;
  this lemma is its injectivity (faithfulness) half rendered as a
  hom-set statement at the level of \(A\)-algebra maps.
\end{proof}

\begin{lemma}\leanok
[Existence of an $A$-algebra map from a $\mathrm{HomOver}$ datum in the \texttt{BijectiveOnStalks} setting --- Archon split]
  \label{lem:identifies-local-ring-spec-existence}
  \lean{Algebra.BijectiveOnStalks.exists_algHom_of_continuousMap}
  \uses{def:identify-local-rings,def:identifies-local-ring-homover,
        lem:identifies-local-ring-spec-spaces-over-X}
  Let \(A\) be a ring and let \(A \to B\), \(A \to C\) be two
  \(A\)-algebra structures both identifying local rings. For every
  \((\varphi, h) \in \mathrm{HomOver}(A, B, C)\) there exists an
  \(A\)-algebra homomorphism \(f \colon B \to C\) with
  \(\Spec(f) = \varphi\) (equivalently:
  \(\Phi(f) = (\varphi, h)\)). This is the substantive content of
  Stacks 096L --- the surjectivity half of the bijection
  \(\Phi\).
\end{lemma}

\begin{proof}
  \leanok
  \uses{def:identify-local-rings,def:identifies-local-ring-homover,
        lem:identifies-local-ring-spec-spaces-over-X,
        lem:identifies-local-ring-invImage-structureSheaf-iso}
  \textit{Source: Stacks Project, Tag 096L (the surjectivity half of
  the cited bijection); the structure-sheaf identification used here
  appears in Tag 096L's proof and uses Tag 096K as its main reduction
  step.}

  Write \(X = \Spec A\), \(Y = \Spec B\), \(Z = \Spec C\), with
  structure maps \(p \colon Y \to X\) and \(q \colon Z \to X\). The
  hypothesis that \(A \to B\) (resp.\ \(A \to C\)) identifies local
  rings is equivalent --- as recorded inside the proof of Stacks 096L
  --- to the structure-sheaf identification
  \(\mathcal O_Y = p^{-1}\mathcal O_X\) (resp.\
  \(\mathcal O_Z = q^{-1}\mathcal O_X\)). The compatibility witness
  \(h\) of \((\varphi, h)\) says exactly \(p \circ \varphi = q\) on the
  topological level.

  Apply \Cref{lem:identifies-local-ring-spec-spaces-over-X} to the
  pair \((p \colon Y \to X, q \colon Z \to X)\). Under
  \(\mathcal O_Y = p^{-1}\mathcal O_X\), the LRS-over-\(X\) forgetful
  map is bijective, so the continuous map \(\varphi \colon Z \to Y\)
  promotes uniquely to a morphism of locally ringed spaces
  \(\tilde\varphi \colon (Z, \mathcal O_Z) \to (Y, \mathcal O_Y)\) over
  \((X, \mathcal O_X)\). Taking global sections of the induced sheaf
  map \(\tilde\varphi^\sharp \colon \mathcal O_Y \to \varphi_* \mathcal O_Z\)
  yields a ring map
  \[
    \mathcal O_Y(Y) \;\longrightarrow\; \mathcal O_Z(Z),
      \qquad \text{i.e.}\qquad B \to C.
  \]
  Call this \(f \colon B \to C\). Tracing the global sections of the
  diagram \(\mathcal O_X(X) \to \mathcal O_Y(Y) \to \mathcal O_Z(Z)\)
  shows \(f \circ \alpha_B = \alpha_C\), so \(f\) is an \(A\)-algebra
  homomorphism. The Spec-map of \(f\) agrees with \(\varphi\) on
  points because the LRS-promotion \(\tilde\varphi\) of \Cref{lem:identifies-local-ring-spec-spaces-over-X} pinpoints the unique LRS-morphism over \(X\) with underlying map \(\varphi\), and \(\Spec(f)\) is another such morphism (with underlying map \(\Spec(f)\) and sheaf part the canonical comap of \(f\)); both must coincide. Hence \(\Spec(f) = \varphi\), as required.

  \smallskip

  \noindent\emph{Warning: the naive per-prime route does not close
  existence.}
  % iter-134 LM-vs-BP-checker `identifies-iter134` noted this trap;
  % the iter-134 mathlib-analogist `lane-d-lrs-bridge` recipe walked
  % into it.
  A tempting strategy --- and the one followed by the iter-134
  prover round --- is to mimic the injectivity argument by building
  \(f(b)\) per-prime: for each \(\mathfrak p \subset C\) compute the
  candidate image of \(b \in B\) in \(C_\mathfrak p\) via the chain
  \(B \to B_{\varphi(\mathfrak p)} \cong A_{\varphi(\mathfrak p) \cap A}
  = A_{\mathfrak p \cap A} \cong C_\mathfrak p\), then assemble these
  into an element of \(\prod_{\mathfrak p} C_\mathfrak p\), and finally
  recover an element of \(C\) via the injection
  \(C \hookrightarrow \prod_{\mathfrak p} C_\mathfrak p\). The first
  three steps go through unchanged from the injectivity proof, but the
  final step is \emph{not} a formality: the injection \(C \hookrightarrow \prod_{\mathfrak p} C_\mathfrak p\)
  is generally not surjective, and showing that the candidate family
  lies in its image is precisely the structure-sheaf identification
  \(\mathcal O_Z = q^{-1}\mathcal O_X\) repackaged as a "local data
  globalises" sheaf condition. Per the iter-134 prover's task report
  (\texttt{.archon/task\_results/Algebra\_IdentifiesLocalRings.lean.md})
  this is the trap: the algebraic-only per-prime route eliminates the
  \texttt{LocallyRingedSpace.Hom} scaffolding but does not, on its
  own, close the existence half. Some form of the structure-sheaf
  identification --- either the LRS route above or one of the
  Mathlib-side alternatives listed next --- must appear in the proof.

  \smallskip

  \noindent\emph{Committed Mathlib route: sheaf-theoretic short circuit (iter-136).}
  The iter-135 \texttt{mathlib-analogist} subagent
  (\texttt{lane-d-sub-existence}) returned an
  \texttt{ALIGN\_WITH\_MATHLIB} verdict on the \emph{sheaf-theoretic
  short circuit} route; the plan agent has committed to this route
  for iter-136. The full recipe lives in
  \texttt{analogies/identifies-local-rings-lrs-bridge.md} (iter-135
  revision). The three Mathlib idioms that close the route are:
  \begin{itemize}
    \item \texttt{AlgebraicGeometry.Spec.locallyRingedSpaceMap} ---
      promotes a ring map to a morphism of locally ringed spaces
      \(\Spec(\text{target}) \to \Spec(\text{source})\), packaging the
      structure-sheaf comap automatically. We use it implicitly via the
      structure-sheaf identification supplied by the helper below: once
      we know \(q^{-1}\mathcal O_Y \cong \mathcal O_Z\) on \(\Spec C\),
      the underlying continuous map \(\varphi \colon Z \to Y\) promotes
      uniquely to an LRS-morphism \(\tilde\varphi \colon (Z,\mathcal O_Z)
      \to (Y,\mathcal O_Y)\) over \((X,\mathcal O_X)\) by
      \Cref{lem:identifies-local-ring-spec-spaces-over-X}.
    \item \texttt{AlgebraicGeometry.Spec.preimage} (or, if the
      one-shot form is unavailable in the current Mathlib snapshot,
      its source \texttt{AlgebraicGeometry.ΓSpec.adjunction.homEquiv})
      --- extracts the underlying ring map \(B \to C\) from an
      LRS-morphism \(\tilde\varphi\) via global sections / \(\Spec\)-\(\Gamma\)
      adjunction.
    \item the new project helper
      \Cref{lem:identifies-local-ring-invImage-structureSheaf-iso}
      (stated below) --- supplies the structure-sheaf identification
      \(q^{-1}\mathcal O_Y \cong \mathcal O_Z\) on \(\Spec C\) that, in
      combination with \Cref{lem:identifies-local-ring-spec-spaces-over-X},
      upgrades the continuous \(\varphi\) to an LRS-morphism over
      \(X\). The helper is proved stalk-locally using
      \texttt{Algebra.BijectiveOnStalks.localRingHom\_comp\_stalkIso}
      (already in \texttt{Proetale/Algebra/StalkIso.lean}) and
      \texttt{Algebra.BijectiveOnStalks.bijective\_localRingHom}.
  \end{itemize}
  Concretely, the existence proof first applies
  \Cref{lem:identifies-local-ring-invImage-structureSheaf-iso} (with
  the data \(A \to C\)) to obtain the structure-sheaf identification
  \(q^{-1}\mathcal O_X \cong \mathcal O_Z\) on \(\Spec C\). Together
  with \Cref{lem:identifies-local-ring-spec-spaces-over-X} applied to
  \((p \colon Y \to X, q \colon Z \to X)\) --- whose hypothesis
  \(\mathcal O_Y = p^{-1}\mathcal O_X\) is supplied by the same helper
  applied to \(A \to B\) --- this upgrades the underlying continuous
  map \(\varphi \colon Z \to Y\) to an LRS-morphism \(\tilde\varphi\)
  over \((X,\mathcal O_X)\). The ring map \(f \colon B \to C\) is then
  extracted via \(\Spec/\Gamma\) (i.e.\ \texttt{Spec.preimage}). The
  compatibility witness \(h\) of \((\varphi, h)\) supplies the
  over-\(X\) structure.

  \smallskip

  \noindent\emph{Deferred routes (iter-136 will not pursue).}
  The two algebraic-only alternatives explored in earlier iters are
  recorded here for completeness but are deliberately \emph{not}
  followed in iter-136:
  \begin{itemize}
    \item \emph{Flat-epimorphism / Olivier characterisation.} An
      algebra map identifying local rings is the same as a flat
      epimorphism (Olivier's theorem); the iter-135 analogist
      estimated this requires a Mathlib gap fill of \(\geq 500\)
      LOC before it could plug in here, so it is deferred.
    \item \emph{Direct sheafification.} An inline sheafification
      argument on basic opens \(D(c) \subset \Spec C\) closes the
      statement but produces no reusable artefact for downstream
      \texttt{BijectiveOnStalks} consumers; deferred for that
      reason.
  \end{itemize}

  \smallskip

  \noindent\emph{Common helpers (all routes).}
  \begin{itemize}
    \item \texttt{PrimeSpectrum.comap\_comp} --- the composition
      functoriality of \(\Spec\) on ring maps. Used in the proof that
      \(\Spec(f) = \varphi\).
    \item \texttt{Localization.localRingHom} and
      \texttt{Algebra.BijectiveOnStalks.localRingHom} (already in
      \texttt{Proetale/Algebra/StalkIso.lean}) --- give the per-prime
      ring isomorphisms \(A_{\mathfrak q \cap A} \cong B_\mathfrak q\)
      and \(A_{\mathfrak p \cap A} \cong C_\mathfrak p\) used in
      every route.
    \item \texttt{PrimeSpectrum.toPiLocalization\_injective} --- the
      injection \(C \hookrightarrow \prod_{\mathfrak p} C_\mathfrak p\).
      Used in the injectivity sub-lemma; in the existence sub-lemma it
      appears as the target object whose image-of-\(C\) characterisation
      is the trap discussed above.
  \end{itemize}
\end{proof}

% NOTE (review iter-142): Lane D-sub prover (iter-142, second consecutive
% iter against this signature) reports the named helper
% `Algebra.BijectiveOnStalks.isIso_specMap_c` is mathematically wrong as
% stated. The c-component of `Spec.locallyRingedSpaceMap` lives on
% `Spec A` (a sheaf morphism `O_{Spec A} -> q_* O_{Spec B}`) and is iso
% iff the ring map is a surjection on both stalks AND global sections,
% strictly stronger than `Algebra.BijectiveOnStalks A B` (counterexample:
% `Z -> Q`; see `Proetale/Algebra/IdentifiesLocalRings.lean` L210-230
% docstring). The corrected helper lives on `Spec B`:
%   `Algebra.BijectiveOnStalks.isIso_invImage_structureSheaf` --- the
%   inverse-image map `q.base^{-1} O_{Spec A} -> O_{Spec B}` is iso.
% This is the Stacks 096J <=> 04D2 sheaf-iso at the inverse-image level.
% Plan agent: please dispatch a blueprint-writer to (a) rename
% `lem:identifies-local-ring-spec-iso-c` to
% `lem:identifies-local-ring-invImage-structureSheaf-iso`, (b) restate
% the lemma at the inverse-image level, (c) update the `\lean{...}` to
% `Algebra.BijectiveOnStalks.isIso_invImage_structureSheaf`. See
% `task_results/Algebra_IdentifiesLocalRings.lean.md` (iter-142) for the
% full closure recipe + Option B (algebraic re-packaging via
% `PrimeSpectrum.toPiLocalization`).
\begin{lemma}\leanok
[Inverse-image of $\mathcal O_{\Spec A}$ is $\mathcal O_{\Spec B}$ when $A \to B$ identifies local rings --- Archon helper]
  \label{lem:identifies-local-ring-invImage-structureSheaf-iso}
  \lean{Algebra.BijectiveOnStalks.isIso_invImage_structureSheaf}
  \uses{def:identify-local-rings, lem:identifies-local-ring-spec-spaces-over-X}
  Let $A \to B$ be a ring homomorphism such that $A \to B$ identifies
  local rings (i.e.\ \texttt{Algebra.BijectiveOnStalks A B}), and let
  $q : \Spec B \to \Spec A$ denote the induced morphism of locally
  ringed spaces. Then the canonical inverse-image map on structure
  sheaves
  \[
    q^{-1}\mathcal O_{\Spec A} \;\longrightarrow\; \mathcal O_{\Spec B}
  \]
  (the morphism of sheaves of rings on $\Spec B$ given by the
  $q^{-1} \dashv q_*$ adjunction counit applied to the structure-sheaf
  comap) is an isomorphism.
\end{lemma}

\begin{proof}
  \uses{def:identify-local-rings}
  \leanok
  The statement is local on $\Spec B$: a morphism of sheaves of rings
  on $\Spec B$ is an isomorphism iff its stalk map at every prime is
  an isomorphism. For each $\mathfrak p \subset B$, the stalk of
  $q^{-1}\mathcal O_{\Spec A}$ at $\mathfrak p$ is canonically
  identified with the stalk of $\mathcal O_{\Spec A}$ at
  $q(\mathfrak p) = \mathfrak p \cap A$ (Mathlib:
  \texttt{TopCat.Presheaf.stalkPullback}), and standard Mathlib
  identifies these structure-sheaf stalks of $\Spec$ with localizations,
  giving
  \[
    (q^{-1}\mathcal O_{\Spec A})_{\mathfrak p}
    \;\cong\; (\mathcal O_{\Spec A})_{\mathfrak p \cap A}
    \;\cong\; A_{\mathfrak p \cap A},
    \qquad
    (\mathcal O_{\Spec B})_{\mathfrak p}
    \;\cong\; B_{\mathfrak p}.
  \]
  Under these identifications the induced stalk map agrees with the
  local ring map
  \texttt{Localization.localRingHom} attached to the data $A \to B$
  (i.e.\ the underlying ring hom witnessed by
  \texttt{RingHom.BijectiveOnStalks.localRingHom} from
  \texttt{Proetale/Algebra/StalkIso.lean}). By
  \texttt{Algebra.BijectiveOnStalks.bijective\_localRingHom} this map
  is bijective, hence an iso of (local) rings.

  \medskip
  \noindent \emph{Sheaf-vs-presheaf stalk identification.} The identification
  $(q^{-1}\mathcal O_{\Spec A})_{\mathfrak p} \cong (\mathcal O_{\Spec A})_{q(\mathfrak p)}$
  above is at the \emph{presheaf} level: Mathlib's
  \texttt{TopCat.Presheaf.stalkPullbackIso} identifies the stalk of the
  \emph{presheaf} pullback with the source-stalk at the image point. The
  inverse-image $q^{-1}\mathcal O_{\Spec A}$ that actually appears in the
  lemma statement is the \emph{sheaf}-level pullback, defined via the
  sheafification adjunction
  $\texttt{TopCat.Sheaf.pullbackPushforwardAdjunction}$ as the sheafification
  of the presheaf pullback:
  \[
    q^{-1}\mathcal O_{\Spec A}
    \;=\; \texttt{sheafify}\bigl(
      \texttt{Presheaf.pullback}_q\,\mathcal O_{\Spec A}^{\mathrm{pre}}
    \bigr).
  \]
  The bridge from sheaf stalk to presheaf stalk is the
  \emph{sheafification-preserves-stalks} principle: for any presheaf
  $F$ of $\mathrm{CommRingCat}$-valued objects on a topological space
  $X$ and any point $x \in X$, the sheafification unit
  $\eta : F \to \texttt{sheafify}\,F$ induces an isomorphism on stalks
  \[
    \eta_x : F_x \;\xrightarrow{\sim}\; (\texttt{sheafify}\,F)_x.
  \]
  For $\mathrm{Type}$-valued presheaves this is Mathlib's
  \texttt{TopCat.Presheaf.sheafifyStalkIso}
  (\texttt{Mathlib/Topology/Sheaves/Sheafify.lean}). For $\mathrm{CommRingCat}$-valued
  presheaves the analogous lemma is \emph{not present in Mathlib}; the file
  \texttt{Sheafify.lean} explicitly carries an
  \texttt{assert\_not\_exists CommRingCat} declaration. This is the
  substantive Mathlib gap blocking the close of the helper.

  The Lean delivery
  combines this with the structure-sheaf-stalk identifications shipped
  by Mathlib's
  \texttt{AlgebraicGeometry.StructureSheaf.stalkIso} and concludes
  via Mathlib's `iso of stalk-iso' principle.
  % NOTE: iter-144 refactor verified that no declaration named
  % `localRingHom_comp_stalkIso` exists in `Proetale/Algebra/StalkIso.lean`;
  % Mathlib's `AlgebraicGeometry.localRingHom_comp_stalkIso` (in
  % `Mathlib/AlgebraicGeometry/Spec.lean`) is the correct carrier and is
  % used directly in the closure. The iter-146 prover round identified the
  % sheafification-preserves-stalks gap for CommRingCat as the substantive
  % remaining obstacle (see Lemma
  % `lem:topcat-presheaf-sheafify-stalk-iso-concrete` below). The iter-147
  % prover task is to formalize that helper lemma in a Mathlib-shaped
  % supplementary file (or inline in StructureSheafPullback.lean) before
  % closing the headline `isIso_invImage_structureSheaf`.
\end{proof}

\begin{lemma}[Sheafification preserves stalks for sheaves valued in a concrete category whose forgetful functor preserves filtered colimits --- Archon Mathlib gap]
  \label{lem:topcat-presheaf-sheafify-stalk-iso-concrete}
  \lean{TopCat.Presheaf.sheafifyStalkIso_concrete}
  Let $\mathcal C$ be a category with all small colimits and limits,
  equipped with a faithful forgetful functor
  $\mathrm{forget}_{\mathcal C} : \mathcal C \to \mathrm{Type}$ that
  preserves filtered colimits, preserves limits, and reflects
  isomorphisms (these are the standard conditions on the forgetful
  functor of a concrete category that occur for $\mathrm{AddCommGrp}$,
  $\mathrm{Ring}$, $\mathrm{CommRingCat}$, $\mathrm{Module}_R$, etc.).
  Let $X : \mathrm{TopCat}$ be a topological space and let
  $F : \mathrm{TopCat.Presheaf}\,\mathcal C\,X$ be a presheaf of
  $\mathcal C$-valued sections on $X$. Let $x \in X$ be a point. Then
  the sheafification unit
  $\eta : F \to \texttt{presheafToSheaf}\,F$ induces an isomorphism on
  stalks: the map
  \[
    \eta_x \;:=\; (\texttt{stalkFunctor}\,\mathcal C\,x).\mathrm{map}(\eta.\mathrm{app}\,F)
    \;:\; F_x \;\longrightarrow\; (\texttt{presheafToSheaf}\,F).\mathrm{val.stalk}\,x
  \]
  is an isomorphism in $\mathcal C$.
  \uses{}
\end{lemma}
\begin{proof}
  Mathlib's \texttt{TopCat.Presheaf.sheafifyStalkIso} establishes the
  same fact for the case $\mathcal C = \mathrm{Type}$. The forgetful
  functor $\mathrm{forget}_{\mathcal C} : \mathcal C \to \mathrm{Type}$
  commutes with the stalk functor (the stalk is a filtered colimit and
  $\mathrm{forget}_{\mathcal C}$ preserves filtered colimits). It also
  commutes with the sheafification functor (sheafification can be
  computed pointwise as a filtered colimit over covers, and
  $\mathrm{forget}_{\mathcal C}$ preserves filtered colimits).
  Therefore the map
  $\mathrm{forget}_{\mathcal C}(\eta_x)$ identifies with the
  $\mathrm{Type}$-level sheafification-stalk iso applied to
  $\mathrm{forget}_{\mathcal C}\,F$, and is hence an iso in
  $\mathrm{Type}$. Since $\mathrm{forget}_{\mathcal C}$ reflects
  isomorphisms, $\eta_x$ is an iso in $\mathcal C$.
\end{proof}

\noindent The \texttt{CommRingCat}-valued instance of
\cref{lem:topcat-presheaf-sheafify-stalk-iso-concrete} is the missing
ingredient in the close of
\cref{lem:identifies-local-ring-invImage-structureSheaf-iso}. Once it
is in scope, the full stalk identification chain becomes
\begin{align*}
  \bigl((\texttt{Sheaf.pullback}\,q.\mathrm{base}).\mathrm{obj}\,\mathcal O_{\Spec A}\bigr).\mathrm{val.stalk}\,\mathfrak p
  & \;\xrightarrow{\texttt{pullbackIso}}
    \bigl(\texttt{sheafify}\,(\texttt{Presheaf.pullback}\,q.\mathrm{base}\,\mathcal O_{\Spec A}^{\mathrm{pre}})\bigr).\mathrm{val.stalk}\,\mathfrak p \\
  & \;\xrightarrow{\texttt{sheafifyStalkIso\_concrete}}
    (\texttt{Presheaf.pullback}\,q.\mathrm{base}\,\mathcal O_{\Spec A}^{\mathrm{pre}}).\mathrm{stalk}\,\mathfrak p \\
  & \;\xrightarrow{\texttt{stalkPullbackIso}}
    \mathcal O_{\Spec A}^{\mathrm{pre}}.\mathrm{stalk}\,(q.\mathrm{base}\,\mathfrak p) \\
  & \;\xrightarrow{\texttt{StructureSheaf.stalkIso}}
    A_{\mathfrak p \cap A} \\
  & \;\xrightarrow{\texttt{localRingHom}}
    B_{\mathfrak p}\quad \text{(bijection from Algebra.BijectiveOnStalks)} \\
  & \;\xrightarrow{\texttt{StructureSheaf.stalkIso.symm}}
    \mathcal O_{\Spec B}.\mathrm{stalk}\,\mathfrak p.
\end{align*}
Compatibility of the actual stalk map of
$q^{-1}\mathcal O_{\Spec A} \to \mathcal O_{\Spec B}$ with this
composite follows from
\texttt{AlgebraicGeometry.localRingHom\_comp\_stalkIso} after unwinding
the adjunction transpose $\texttt{homEquiv.symm}$ of the direct-image
$c$-component $q.\mathrm{toHom.c}$.
% ------------------------------------------------------------------
% End iter-135 split sub-lemmas.
% ------------------------------------------------------------------

\begin{definition}
  \label{def:identifies-local-ring-hom-set-bijection}
  \lean{Algebra.BijectiveOnStalks.algHomEquivContinuousMap,
        Algebra.BijectiveOnStalks.algHom_of_continuousMap,
        Algebra.BijectiveOnStalks.continuousMap_of_algHom}
  \uses{def:identify-local-rings,def:identifies-local-ring-homover,thm:identifies-local-ring-to-top-fully-faithful}
  Let \(A\) be a ring and \(X = \Spec(A)\). Let \(A \to B\) and \(A \to C\) be two \(A\)-algebras that identifies local rings. As a consequence of \Cref{thm:identifies-local-ring-to-top-fully-faithful}, we have a bijective map
  \[F : \Hom_{A\textnormal{-alg}}(B, C) \to \mathrm{HomOver}(A, B, C),\]
  with \(\mathrm{HomOver}(A, B, C)\) as introduced in
  \Cref{def:identifies-local-ring-homover}, sending
  \(f : B \to C\) to the continuous map
  \(\Spec(f) : \Spec(C) \to \Spec(B)\) induced by \(f\) (paired with
  its automatic compatibility witness over \(\Spec(A)\)).
\end{definition}

\begin{lemma}
    Let $A = \colim_i A_i$ be a filtered colimit of $A$-algebras and let $B$ be an étale $A$-algebra. Then
    there exists an $i$ and an étale $A_i$-algebra $B'$ such that
    \[
    \begin{tikzcd}
        A_i \arrow{r} \arrow{d} & B' \arrow{d} \\
        A \arrow{r} & B
    \end{tikzcd}
    \] is a pushout diagram.
    \label{lemma:etale-ind-spreads}
\end{lemma}

\begin{proof}
    This follows because every étale algebra is of finite presentation and by possibly enlarging $i$ we can ensure
    that $B'$ is étale over $A_i$.
\end{proof}

\section{ind-Zariski maps}

\begin{definition}[Ind-Zariski, {\cite[\href{https://stacks.math.columbia.edu/tag/096N}{Tag 096N}]{stacks-project}}]
    \label{def:ind-Zariski}
    \uses{def:local-isomorphism}
    \lean{Algebra.IndZariski}
    \leanok
    An $R$-algebra $S$ is called a \emph{ind-Zariski} if $S$ can be written as a filtered colimit $S \simeq \colim S_i$ with each $R \to S_i$ a local isomorphism. A ring homomorphism $f : R \to S$ is called a \emph{ind-Zariski} if $S$ is ind-Zariski as an $R$-algebra via $f$.
\end{definition}

\begin{remark}
  The definition of ind-Zariski is slightly more general from \cite[Definition 2.2.1(iv)]{proetale}, where they defined as a filtered colimit of finite products of principal localizations.
\end{remark}

\begin{lemma}[{\cite[\href{https://stacks.math.columbia.edu/tag/096T}{Tag 096T}]{stacks-project}}]
  Let $A \to B$ be an ind-Zariski ring map. Then it identifies local rings, i.e. for every prime 
  $\mathfrak{q} \subset B$ the canonical map $A_{\varphi^{-1}(\mathfrak{q})} \to B_{\mathfrak{q}}$
  is an isomorphism.
  \label{thm:ind-Zariski-identifies-local-rings}
  \uses{def:ind-Zariski,def:identify-local-rings}
  \lean{Algebra.IndZariski.bijectiveOnStalks_algebraMap,RingHom.IndZariski.bijectiveOnStalks}
\end{lemma}

\begin{proof}
  Omitted.
\end{proof}

\begin{lemma}
  Let $A \to B$ be an ind-Zariski ring map. Then $A \to B$ is flat.
  \label{thm:ind-Zariski-is-flat}
  \uses{def:ind-Zariski}
  \lean{Module.Flat.of_indZariski,RingHom.IndZariski.flat}
\end{lemma}

\begin{proof}
    An ind-Zariski ring map is a filtered colimit of local isomorphisms. A local isomorphism is flat.
\end{proof}

\begin{lemma}
  A filtered colimit of ind-Zariski algebras is ind-Zariski.
  \label{lemma:ind-ind-Zariski}
  \uses{def:ind-Zariski}
  \lean{Algebra.IndZariski.of_colimitPresentation,RingHom.IndZariski.of_isColimit,
    Algebra.IndZariski.iff_ind_indZariksi,RingHom.IndZariski.iff_ind_indZariski}
\end{lemma}

\begin{proof}
    This follows from general theory, because local isomorphisms are of finite presentation.
\end{proof}

\begin{lemma}
  A finite product of ind-Zariski algebras is ind-Zariski.
  \label{lemma:ind-Zariski-products}
  \uses{def:ind-Zariski}
  \lean{Algebra.IndZariski.prod,
        RingHom.IndZariski.prod,
        Algebra.IndZariski.pi,
        RingHom.IndZariski.pi}
\end{lemma}

\begin{proof}
  This is the consequence of the more general fact that if a property $p$ of ring maps is stable
  under finite products, then Ind-$p$ is also stable under finite products since finite product
  of filtered index category is again filtered.

  It suffices to show that being local isomorphism is stable under finite products. Suppose
  $A \to B_i$ is a local isomorphism for each $i$ for finitely many $i$'s. Every prime $\q$ in
  $\prod_i B_i$ is an extension prime ideal of one of the factors, say $B_0$. Since there exists
  $f_0 \in B_0$, $f_0 \notin \q$ such that $A \to (B_0)_{f_0}$ induces an open immersion,
  localize $\prod_i B_i$ at the corresponding idempotent times this element \(f_0\) to obtain
  the desired result.
\end{proof}

\begin{lemma}
  \label{thm:ind-Zariski-composition}
  Let \(A \to B\) and \(B \to C\) be ind-Zariski ring maps. Then the composition \(A \to C\) is also ind-Zariski.
  \uses{def:ind-Zariski}
  \lean{Algebra.IndZariski.trans,RingHom.IndZariski.comp}
\end{lemma}

\begin{proof}
  Omitted.
\end{proof}

\begin{lemma}
    \label{lemma:ind-Zariski-localization}
    Let $A$ be a ring and $S \subseteq A$ a multiplicative subset. Then $S^{-1}A$ is ind-Zariski over $A$.
    \lean{Algebra.IndZariski.of_isLocalization}
    \leanok
\end{lemma}

\begin{proof}
    \leanok
    $S^{-1} A = \colim_{f \in S} A_f$ and $A \to A_f$ is a local isomorphism.
\end{proof}

\section{Henselian objects and Henselisation}

In this section we define a generalisation of the notion of a Henselian ring to an arbitrary category. Let $\mathcal{C}$ be a category
and $P$ a property of morphisms.

\begin{definition}[Henselian]
    A morphism $f\colon X \to Y$ is $P$-\emph{Henselian}, if for every factorisation
    \[
    \begin{tikzcd}
        X \arrow{r}{u} \arrow[swap]{dr}{f} & Z \arrow{d}{v} \\
                    & Y,
    \end{tikzcd}
    \] with $u$ satisfying $P$, there exists a retraction of $u$.
    \label{def:p-henselian}
\end{definition}

\begin{example}
    A local ring $(R, \mathfrak{m}, \kappa)$ is Henselian if and only if the projection $R \to \kappa$ is étale-Henselian.
\end{example}

\begin{proof}
    By {\cite[\href{https://stacks.math.columbia.edu/tag/04GG}{Tag 04GG}]{stacks-project}}, $R$ is Henselian if and only if
    for any étale ring map $R \to S$ and prime $\mathfrak{q}$ of $S$ lying over $\mathfrak{m}$ with $\kappa = \kappa(\mathfrak{q})$,
    there exists a retraction $S \to R$ of $R \to S$.

    Suppose $R$ is Henselian and let $R \to S \to \kappa$ be a factorisation of $R \to \kappa$ with $R \to S$ étale. Then
    $\mathfrak{q} = \mathrm{ker}(S \to \kappa)$ is a prime ideal of $S$ lying above $\mathfrak{m}$. Hence we obtain a commutative
    diagram
    \[
        \begin{tikzcd}
        R \arrow{r} \arrow{d} & S \arrow{r} \arrow{d} & \kappa \\
        \kappa \arrow[dashed]{r} & \kappa(\mathfrak{q}) \arrow[dashed]{ur},
        \end{tikzcd}
    \] where the composition of the dashed arrows is the identity of $\kappa$. Hence $\kappa(\mathfrak{q}) = \kappa$ and by assumption,
    $R \to S$ has a retraction.

    Now assume that $R \to \kappa$ is étale-Henselian and let $R \to S$ be étale and $\mathfrak{q}$ be a prime ideal of $S$ with
    $\kappa(\mathfrak{q}) = \kappa$. Hence the composition $R \to S \to \kappa(\mathfrak{q}) = \kappa$ is a factorisation
    of $R \to \kappa$. Thus $R \to S$ has a retraction.
\end{proof}

% TODO: add definition of Henselisation if needed

\begin{definition}[Strictly Henselian local ring]
    \uses{def:p-henselian}
    \label{def:strictly-henselian-local-ring}
    A local ring $(R, \mathfrak{m}, \kappa)$ is \emph{strictly henselian} if the composition $R \to \kappa \to \kappa^{\mathrm{sep}}$ is étale-Henselian
    for a separable closure $\kappa^{\mathrm{sep}}$.
\end{definition}

The following example shows that our definition agrees with the one in
\cite[\href{https://stacks.math.columbia.edu/tag/04GF}{Tag 04GF}]{stacks-project}.

\begin{lemma}
    \label{lem:sh-iff-henselian-and-sep-closed}
    A local ring $(R, \mathfrak{m}, \kappa)$ is strictly Henselian if and only if it is Henselian and $\kappa$ is separably closed.
    \uses{def:strictly-henselian-local-ring}
    \lean{isStrictlyHenselianLocalRing_iff_henselianLocalRing_and_isSepClosed}
    \leanok
\end{lemma}

\begin{proof}
    TBA.
\end{proof}

\begin{lemma}
    \label{lem:etale-locally-standard-etale-at-max}
    \lean{Algebra.Etale.isStandardEtale_localizationAway_of_etale_at_maxIdeal}
    Let \(A\) be a Henselian local ring and \(A \to B\) an étale ring map. Let \(\n \subset B\) be a maximal ideal lying over the maximal ideal \(\m\) of \(A\). Then there exists \(b \in B \setminus \n\) such that the localization \(B_b = B[1/b]\) is a standard étale \(A\)-algebra.
\end{lemma}

\begin{proposition}
    \label{thm:strictly-henselian-good-retraction}
    \uses{def:strictly-henselian-local-ring, lem:etale-locally-standard-etale-at-max}
    Let \(A, \m\) be a strictly henselian local ring. Let \(f : A \to B\) a \'etale ring map with \(\n\) a maximal ideal of \(B\) lying over \(\m\). Then there exists a retraction \(s : B \to A\) of \(A \to B\) such that \(\n = s^{-1}(\m)\).
\end{proposition}

\begin{proof}
  Since \(B/\m B\) is an étale algebra over the residue field \(k = A/\m\), it splits as \(B/\m B \cong k_1 \times \cdots \times k_n\) for finitely many finite separable extensions \(k_i\) of \(k\). Thus \(B/\n\), as a quotient of \(B/\m B\), is also a finite separable extension of \(k\). Choose an inclusion \(B/\n \hookrightarrow k^{\sep}\). Then we have a commutative diagram
  \[
  \begin{tikzcd}
    A \arrow{r} \arrow{d} & B \arrow{d}\arrow{rd} & \\
    k \arrow{r} & B/\n \arrow{r} & k^{\sep}.
  \end{tikzcd}
  \]
  Since \(A \to k^{\sep}\) is étale-Henselian, there exists a retraction \(r : B \to A\) of \(A \to B\). But this retraction may not satisfy \(\n = r^{-1}(\m)\). 
  
  To solve this, notice that there are only finitely many prime ideals \(\n_0 = \n, \n_1, \cdots, \n_k\) of \(B\) lying over \(\m\) and they are all maximal ideals. (directly by the fact that \(B/\m B \cong k_1 \times \cdots \times k_n\)). Since \(\bigcap_{i \ne 0} \n_i \nsubseteq \n\) (otherwise one of \(\n_i \subseteq \n\) which is impossible), we can find element \(b \in \n_i\) for every \(i \ne 0\) that does not fall in $\n$. We may replace \((B, \n)\) by \((B_b, \n B_b)\), then \(B\) is still étale over \(A\) and \(\n\) is the only prime ideal of \(B\) lying over \(\m\). Then the section \(s_b : B_b \to A\) given by the previous paragraph satisfies \(\n B_b = s_b^{-1}(\m)\). Composition with the localization map \(B \to B_b\) gives the desired section \(s : B \to A\).
\end{proof}

\begin{proposition}
    \label{thm:etale-over-strictly-henselian-localization-isom}
    \uses{def:strictly-henselian-local-ring}
    Let \(A\) be a strictly Henselian local ring and \(f : A \to B\) an étale ring map. Suppose \(\n\) is a maximal ideal of \(B\) lying over the maximal ideal \(\m\) of \(A\). Then the induced map \(A \to B_\n\) is an isomorphism.
\end{proposition}

\begin{proof}
  \uses{thm:strictly-henselian-good-retraction}
  Let \(s: B \to A\) be the section given by \Cref{thm:strictly-henselian-good-retraction} for the pair \((B, \n)\). Localizing at \(\m\), we obtain the following diagram:
  \[
  \begin{tikzcd}
    A = A_\m \arrow[r, "f_\m"]\arrow[dr, "f'"'] & B_\m \arrow[d] \arrow[r, "s_\m"] & A \\
    & B_\n \arrow[ur, dotted, "s'"'] &
  \end{tikzcd}
  \]
  The dotted arrow \(s'\) exists by the universal property of localization, since \(\n = s^{-1}(\m)\). Moreover, \(s'\) is a section of \(f'\).

  Since \(f_\m\) is étale and \(s_\m\) is a section of \(f_\m\), it follows that \(s_\m\) is also étale by the cancellation property of étale maps. Thus \(s'\), being a composition of a localization and an étale map, is flat. As a flat local ring homomorphism, \(s'\) is faithfully flat and hence injective. Since \(s'\) is also a section of \(f'\), it is surjective. Therefore, \(s'\) is an isomorphism, and the composition \(A \to B_\n\) is an isomorphism.
\end{proof}

\section{Ind-etale ring maps}

\begin{definition}[Ind-étale algebras, {\cite[\href{https://stacks.math.columbia.edu/tag/097I}{Tag 097I}]{stacks-project}}]
    An $R$-algebra $S$ is \emph{ind-étale} if it is a filtered colimit of étale $R$-algebras.
    We say a ring homomorphism $R \to S$ is \emph{ind-étale} if $S$ is ind-étale as an $R$-algebra via $f$.
    \label{def:ind-etale}
    \lean{Algebra.IndEtale, RingHom.IndEtale}
\end{definition}

\begin{lemma}[{\cite[\href{https://stacks.math.columbia.edu/tag/0BSI}{Tag 0BSI}]{stacks-project}}]
  \label{thm:ind-etale-comp}
  \lean{Algebra.IndEtale.trans, RingHom.IndEtale.comp}
  \uses{def:ind-etale}
  Let $A \to B$ and $B \to C$ be ring maps. If $A \to B$ and $B \to C$ are ind-étale,
  then $A \to C$ is ind-étale.
\end{lemma}

\begin{proof}
  Omitted.
\end{proof}

\begin{lemma}[{\cite[\href{https://stacks.math.columbia.edu/tag/0BSH}{Tag 0BSH}]{stacks-project}}]
  \label{thm:ind-etale-base-change}
  \uses{def:ind-etale}
  \lean{RingHom.IndEtale.isStableUnderBaseChange}
  Let $A \to B$ be an ind-étale ring map. For any ring map $A \to C$, the base change 
  $C \to B \otimes_A C$ is ind-étale.
\end{lemma}

\begin{proof}
  Omitted.
\end{proof}

\begin{lemma}
    Let $B = \colim_i B_i$ be a filtered colimit of ind-étale $A$-algebras. Then $B$ is
    an ind-étale $A$-algebra.
    \uses{def:ind-etale}
    \label{lemma:ind-ind-etale}
    \lean{Algebra.IndEtale.of_colimitPresentation,RingHom.IndEtale.iff_ind_indEtale,
      RingHom.IndEtale.of_isColimit,Algebra.IndEtale.iff_ind_indEtale}
\end{lemma}

\begin{proof}
    This follows from general theory, because étale maps are of finite presentation.
\end{proof}

\begin{lemma}
  Let $A \to B$ be an ind-Zariski ring map. Then $A \to B$ is ind-étale.
  \label{thm:ind-Zariski-is-ind-etale}
  \uses{def:ind-Zariski,def:ind-etale}
  \lean{Algebra.IndEtale.of_indZariski,RingHom.IndZariski.indEtale}
\end{lemma}

\begin{proof}
  An ind-Zariski ring map is a filtered colimit of local isomorphisms. A local isomorphism is étale.
\end{proof}

\section{w-local spaces}

\begin{definition}[w-local spaces, {\cite[Definition 2.1.1]{proetale}}, {\cite[\href{https://stacks.math.columbia.edu/tag/096A}{Tag 096A}]{stacks-project}}]
    \label{def:w-local-space}
    \lean{WLocalSpace}
    \leanok
    A topological space \(X\) is \emph{w-local} if it satisfies:
    \begin{enumerate}
        \item \(X\) is spectral,
        \item every point specializes to a unique closed point,
        \item the subspace \(X^c\) of closed points is closed.
    \end{enumerate}
\end{definition}

\begin{definition}[w-local morphisms, {\cite[Definition 2.1.1]{proetale}}, {\cite[\href{https://stacks.math.columbia.edu/tag/096A}{Tag 096A}]{stacks-project}}]
    Let \(X\) and \(Y\) be w-local spaces. A morphism \(f: X \to Y\) is \emph{w-local} if it is spectral and the image of closed points \(f(X^c) \subseteq Y^c\).
    \label{def:w-local-space-map}
    \uses{def:w-local-space}
    \lean{IsWLocalMap}
    \leanok
\end{definition}

\begin{definition}[w-purely-local morphisms]
    Let \(X\) and \(Y\) be w-local spaces. A morphism $f\colon X \to Y$ is \emph{w-purely-local} if
    it $f^{-1}(X^c) = Y^c$.
    \label{def:w-purely-local-space-map}
\end{definition}

Obviously, every w-purely-local morphism is w-local.

\begin{lemma}
    Let \(f: X \to Y\) and \(g: Y \to Z\) be w-local maps between w-local spaces. Then the composition \(g \circ f: X \to Z\) is w-local.
    \label{thm:w-local-map-composition}
    \uses{def:w-local-space, def:w-local-space-map}
    \lean{IsWLocalMap.comp}
    \leanok
\end{lemma}

\begin{proof}
    \leanok
    Omitted.
\end{proof}

\begin{lemma}
    Let $X$ be w-local and $Z \subseteq X$ a subset closed under specialization. Then the generalization
    $Z^g$ in $X$ is still closed under specialization.
    \uses{def:generalization}
    \label{lemma:generalization-specialization-closed}
    \lean{StableUnderSpecialization.generalizationHull_of_wLocalSpace}
    \leanok
\end{lemma}

\begin{proof}
    \leanok
    Let $x \rightsquigarrow y$ in $X$ with $x \in Z^g$. Then $x$ specializes to some $z \in Z$, which
    in turn specializes to a unique closed point $c$ in $X$ because $X$ is w-local.
    Because $Z$ is stable under specialization, $c \in Z$. Also $y$ specializes to a unique closed
    point $c'$ in $X$ and by uniqueness we have $c' = c \in Z$. Hence $y \in Z^g$.
\end{proof}

\begin{lemma}
    Let $X$ be w-local and $Z \subseteq X$ a subset closed under specialization. If
    $Z$ is spectral, it is w-local and the inclusion $Z \to X$ maps closed points to closed points.
    \uses{def:w-local-space}
    \label{lemma:w-local-specialization-closed}
    \lean{Topology.IsEmbedding.wLocalSpace_of_stableUnderSpecialization_range}
    \leanok
\end{lemma}

\begin{proof}
    \leanok
    \uses{lemma:embedding-specializations}
    Since $Z$ is closed under specializations, we have $Z^c = X^c \cap Z$. In particular
    $Z^c$ is closed in $Z$ and $Z \to X$ preserves closed points.
    By \ref{lemma:embedding-specializations}, specializations
    in $Z$ are the same as specializations in $X$. Since $Z$ is closed
    under specializations, every point in $Z$ specializes to a unique
    closed point in $X$ and therefore in $Z$.
\end{proof}

\begin{lemma}[{\cite[\href{https://stacks.math.columbia.edu/tag/096B}{Tag 096B}]{stacks-project}}]
    \label{thm:closed-subspace-w-local}
    \uses{def:w-local-space,def:w-local-space-map}
    Let \(X\) be a w-local space. Then every closed subspace \(Y \subseteq X\) is w-local.
    Moreover, the inclusion map \(Y \to X\) is w-local.
    \lean{Topology.IsClosedEmbedding.wLocalSpace}
    \leanok
\end{lemma}

\begin{proof}
    \leanok
    \uses{lemma:w-local-specialization-closed}
    This follows from \ref{lemma:w-local-specialization-closed}, because
    closed subsets are stable under specialization. The inclusion $Y \to X$ is
    spectral as a closed embedding and hence w-local by the lemma.
\end{proof}

\begin{lemma}
    Let $X$ be a profinite topological space and $x, y \in X$ be distinct. Then
    there exist $U, V \subseteq X$ open and closed such that $X = U \coprod V$ and
    $x \in U$, $y \in V$.
    \label{lemma:profinite-disj-clopen-separation}
    \lean{instTotallySeparatedSpaceOfTotallyDisconnectedSpace}
    \mathlibok
\end{lemma}

\begin{proof}
  \mathlibok
  Write $X = \lim_{i \in I} X_i$ as a directed limit of finite discrete spaces. Let $\pi_i \colon X \to X_i$ be the
  projection maps. A base of the open sets of $X$ is given by the preimages $\pi_i^{-1}(\{z\})$ for $z \in X_i$ and $i \in I$. Every open set in this base is also closed, because its complement is a finite union of such sets.

  Choose an open subset in this base containing $x$ but not $y$, say $U = \pi_{i_0}^{-1}(\{z_0\})$ for some $i_0 \in I$ and $z_0 \in X_{i_0}$. Then $V = X - U$ is also open and closed, and we have $X = U \coprod V$ with $x \in U$ and $y \in V$.   
\end{proof}

\begin{lemma}[{\cite[\href{https://stacks.math.columbia.edu/tag/0905}{Tag 0905}, (6) $\implies$ (1)]{stacks-project}}]
  Let $X$ be a spectral space such that every point is closed. Then $X$ is profinite.
  \label{lemma:spectral-closed-points-profinite}
  \lean{SpectralSpace.totallySeparatedSpace}
  \leanok
\end{lemma}

\begin{proof}
  \leanok
  \uses{lemma:spectral-constructible-closed-closed}
  Since every point is closed, there are no nontrivial specializations between points. By \Cref{item:specialization-closed-spectral-constructible-closed-closed} in \Cref{lemma:spectral-constructible-closed-closed}, all constructible subsets of $X$ are closed. In particular, every compact open subset is also closed. Since $X$ is spectral, thus T0, for any two points there is a quasi-compact open (also closed) containing exactly one of them. Thus $X$ is profinite by \verb`isTotallyDisconnected_of_isClopen_set`.
\end{proof}

\begin{lemma}
    Let $X$ be a w-local space. Then the set of closed points $X^c$ of $X$ is profinite.
    \uses{def:w-local-space}
    \label{lemma:w-local-closed-points-profinite}
    \lean{WLocalSpace.CompactSpace_closedPoints}
    \leanok
\end{lemma}

\begin{proof}
    \leanok
    \uses{lemma:spectral-closed-points-profinite}
    It is a closed subspace of a spectral space, hence spectral. Therefore it is profinite by
    \ref{lemma:spectral-closed-points-profinite}
\end{proof}

\begin{lemma}[{\cite[\href{https://stacks.math.columbia.edu/tag/0A31}{Tag 0A31}, (3) $\implies$ (4)]{stacks-project}}]
  Let $X$ be a spectral space and $E \subseteq X$ a quasi-compact subset. Then the generalization $E^g$ is an intersection of quasi-compact open subsets.
  \label{lemma:spectral-generalization-intersection-open}
  \uses{def:generalization}
  \lean{generalizationHull.eq_sInter_of_isCompact}
  \leanok
\end{lemma}

\begin{proof}
  \leanok
  Let $X$ be a spectral space and let $W$ denote $E^g$. Since every point of $W$ specializes to a point of $E$ we see that an open of $W$ which contains $E$ is equal to $W$. Hence since $E$ is quasi-compact, so is $W$. If $x \in X$, $x \notin W$, then $Z = \overline{\{x\}}$ is disjoint from $W$. Since $W$ is quasi-compact we can find a quasi-compact open $U$ (since there is a open basis of quasi-compact opens of $X$ and finitely many of them would suffice to cover $W$) with $W \subset U$ and $U \cap Z = \emptyset$. We conclude that $W$ is an intersection of quasi-compact open subsets.
\end{proof}

\begin{definition}[Constructible topology, {\stacksproject{08YF}}]
    Let $X$ be a topological space. The \emph{constructible topology} on $X$
    is the topology with subbase of opens the sets $U$ and $X - U$ for
    all quasi-compact opens $U$ of $X$.
    \label{def:constructible-topology}
    \lean{constructibleTopology}
\end{definition}

In other words, a subset of $X$ is open in the constructible topology if and only if
it is a union of finite intersections of sets of the form $U$ or $X - U$ for
a quasi-compact open $U$.

\begin{lemma}[\stacksproject{0901}]
    Let $X$ be a spectral space. Then $X$ is quasi-compact in the constructible topology.
    \uses{def:constructible-topology}
    \label{lemma:spectral-constructible-quasi-compact}
    \lean{compactSpace_withConstructibleTopology}
\end{lemma}

\begin{proof}
    Let $\mathfrak{B}$ be the set of subsets of $X$ that are quasi-compact open or
    the complement of a quasi-compact open.
    By definition of the constructible topology and the Alexander subbasis theorem,
    it suffices to show that every covering $X = \bigcup_{i \in  I} B_i$
    with $B_i \in \mathfrak{B}$ has a finite subcover. Since
    for every $B \in \mathfrak{B}$ also $X - B \in \mathfrak{B}$,
    it suffices to show: If $\{B_i\}_{i \in I}$ is a family of elements
    of $\mathfrak{B}$ such that all finite subfamilies have non-empty intersection,
    then $\bigcap_{i \in  I} B_i$ is non-empty.

    Let $\mathcal{S}$ be the set of families in $\mathfrak{B}$ that are pairwise unequal, have non-empty
    finite intersections and empty intersection. For the sake of contradiction, suppose
    $\mathcal{S}$ is non-empty. If we equip $\mathcal{S}$ with the ordering by subset inclusion,
    Zorn's Lemma yields a maximal family $\{B_i\}_{i \in I} \in \mathcal{S}$.

    Let $I' \subseteq I$ be the indices such that $B_i$ is closed and set $Z = \bigcap_{i \in  I'} B_i$.
    This is a closed and non-empty subset of $X$, because $X$ is quasi-compact
    and all finite intersections are non-empty by assumption.
    Note that for any $i_1, \ldots, i_n \in I$ the intersection $\bigcap_{a=1}^n B_{i_a}$ is quasi-compact,
    hence $Z \cap \bigcap_{a=1}^n B_{i_a}$ is non-empty by the same argument.

    Suppose $Z = Z' \cup Z'$ is reducible with $Z', Z'' \subsetneq Z$ closed. Because
    $X$ is a spectral space, there exist quasi-compact opens $U', U'' \subseteq X$ with
    $U' \cap Z'$ non-empty, $U'' \cap Z''$ non-empty and $U' \cap Z'' = \emptyset = U'' \cap Z'$.
    Then set $B' = X - U' \in \mathfrak{B}$ and $B'' = X - U'' \in \mathfrak{B}$. We
    claim that $\{B_i\}_{i \in I}$ with $B'$ or $B''$ added is still in $\mathcal{S}$. Suppose,
    this is not the case. Then there exist finitely many indices
    $i_a$ and $j_b$ in $I$ such that
    \[
    B' \cap \bigcap_{a=1}^n B_{i_a} = \emptyset = B'' \cap \bigcap_{b=1}^m B_{j_b}
    .\] By construction, $Z \subseteq B' \cup B''$, so in particular
    \[
    Z \cap \bigcap_{a=1}^n B_{i_a} \cap \bigcap_{b=1}^m B_{j_b} = \emptyset
    .\] which is a contradiction. Hence $\{B_i\}_{i \in I}$ is not maximal, also a contradiction. So
    $Z$ is irreducible. But for every $i \in I - I'$, the set $B_i$ is open and hence
    the non-empty open $Z \cap B_i$ of $Z$ contains the unique generic point of $Z$. Hence
    $\bigcap_{i \in  I} B_i = Z \cap \bigcap_{i \in  I - I'} B_i$ is non-empty, which is again a contradiction.
\end{proof}

\begin{lemma}[\stacksproject{0903}]
    \uses{def:constructible-topology}
    \label{lemma:spectral-constructible-closed-closed}
    Let $X$ be a spectral space and $E \subseteq X$ a subset closed in the constructible topology.
    \begin{enumerate}
        \item If $x \in \overline{E}$, then $x$ is the specialization of a point in $E$.
        \item If $E$ is stable under specialization, then $E$ is closed. \label{item:specialization-closed-spectral-constructible-closed-closed}
    \end{enumerate}
    \lean{exists_specializes_of_isClosed_constructibleTopology_of_mem_closure, IsClosed.of_isClosed_constructibleTopology}
\end{lemma}

\begin{proof}
    \uses{lemma:spectral-constructible-quasi-compact}
    It suffices to show the first claim. Let $x \in \overline{E}$
    and let $\{U_i\}_{i \in I}$ be the set of quasi-compact open neighbourhoods of $x$.
    Since $X$ is spectral, these form a neighbourhood basis of $x$, so $\bigcap_{i \in  I} U_i$
    is the set of generalizations of $x$. It hence suffices to show
    that $\bigcap_{i \in  I} U_i \cap E$ is non-empty. Indeed,
    as $X$ is quasi-separated, any finite intersection of the $U_i$ is another one and
    since $x \in \overline{E}$, the intersection $U_i \cap E$ is non-empty for all $i$.
    Since $U_i \cap E$ is closed in the constructible topology for all $i$
    and $X$ is quasi-compact in the constructible topology by \ref{lemma:spectral-constructible-quasi-compact},
    the intersection $\bigcap_{i \in  I} U_i \cap E$ is non-empty.
\end{proof}

\begin{lemma}
    Let $X$ be a w-local space and $Z \subseteq X$ a closed subset. Then $Z^g$ is closed.
    \uses{def:w-local-space, def:generalization}
    \label{lemma:w-local-generalization-closed-closed}
    \lean{isClosed_generalizationHull_of_wLocalSpace}
    \leanok
\end{lemma}

\begin{proof}
    \leanok
    \uses{lemma:spectral-generalization-intersection-open, lemma:generalization-specialization-closed, lemma:spectral-constructible-closed-closed}
    Since $Z$ is closed, it is quasi-compact. Hence by \Cref{lemma:spectral-generalization-intersection-open},
    we see that $Z^g$ is closed in the constructible topology as an intersection of quasi-compact opens.
    Since $Z$ is stable under specialization, the same holds for $Z^g$ by
    \Cref{lemma:generalization-specialization-closed}. Thus
    by \Cref{lemma:spectral-constructible-closed-closed}, the claim follows.
\end{proof}

\begin{lemma}[{\cite[\href{https://stacks.math.columbia.edu/tag/0968}{Tag 0968}]{stacks-project}}]
    \label{thm:closed-points-isom-pi0}
    \uses{def:w-local-space}
    Let \(X\) be a w-local space. Then the composition $X^c \to X \to \pi_0(X)$ is a homeomorphism,
    where $X^c$ denotes the closed points of $X$.
    \lean{WLocalSpace.isHomeomorph_connectedComponents_closedPoints}
    \leanok
\end{lemma}

\begin{proof}
    \leanok
    \uses{thm:pi0-profinite, lemma:w-local-closed-points-profinite, lemma:profinite-disj-clopen-separation, def:generalization, lemma:w-local-generalization-closed-closed}
    Since $X^c$ is quasi-compact and $\pi_0(X)$ is Hausdorff by \Cref{thm:pi0-profinite}, it suffices
    to show the map is bijective.
    To show injectivity let $x, y \in X^c$ be distinct. Since $X^c$ is profinite
    (\Cref{lemma:w-local-closed-points-profinite}), there
    exist open and closed subsets $U_0, V_0 \subseteq X^c$ such that
    $X^c = U_0 \coprod V_0$ by \Cref{lemma:profinite-disj-clopen-separation}.
    Let $U = U_0^g$ and $V = V_0^g$. Because $X$ is w-local, we get
    $X = U \coprod V$ as sets. By \Cref{lemma:w-local-generalization-closed-closed},
    both $U$ and $V$ are closed and therefore also both open. Hence
    $x$ and $y$ lie in distinct connected components of $X$, so the images in $\pi_0(X)$ are distinct.

    Since every connected component is closed, it is quasi-compact and hence contains a closed point. Thus
    the map is surjective.
\end{proof}

\begin{lemma}[{\cite[\href{https://stacks.math.columbia.edu/tag/096C}{Tag 096C}, second part]{stacks-project}}]
\label{thm:cartesian-profinite-w-local}
\uses{def:w-local-space,def:w-local-space-map}
\lean{ConnectedComponents.wlocalSpace_of_isPullback,ConnectedComponents.isWLocalMap_of_isPullback,ConnectedComponents.preimage_closedPoints_eq_closedPoints_of_isPullback}
Let $X$ be a spectral space. Let
\[
\begin{tikzcd}
Y \arrow[r] \arrow[d] & T \arrow[d] \\
X \arrow[r] & \pi_0(X)
\end{tikzcd}
\]
be a cartesian diagram in the category of topological spaces with $T$ profinite. If moreover $X$ is w-local, then $Y$ is w-local, $Y \to X$ is w-local, and the set of closed points of $Y$ is the inverse image of the set of closed points of $X$.
\end{lemma}

\begin{proof}
  \uses{thm:pullback-spectral,thm:cartesian-profinite,thm:closed-points-isom-pi0}
  We know that $Y$ is spectral and $T = \pi_0(Y)$ by \Cref{thm:cartesian-profinite}. Assume $X$ is w-local and let $X^c \subset X$ be its closed points. The inverse image $Y^c \subset Y$ of $X^c$ maps bijectively onto $T$, since $X^c \to \pi_0(X)$ is a bijection by \Cref{thm:closed-points-isom-pi0}. (Note that we do not yet know whether $Y^c$ is the set of closed points of $Y$.) Moreover, $Y^c$ is quasi-compact as a closed subset of the spectral space $Y$. Hence $Y^c \to \pi_0(Y) = T$ is a homeomorphism by \verb`isHomeomorph_iff_continuous_bijective` (\stacksproject{08YE}). It follows that all points of $Y^c$ are closed in $Y$: if some $y \in Y^c$ specialized to a distinct closed point $y'$ (which also falls in $Y$ since $Y$ is closed), then both would map to the same point in $T$, contradicting the bijectivity of $Y^c \to T$.

  Conversely, if $y \in Y$ is a closed point, then it is closed in its fibre over $T$, and its image $x \in X$ is closed in the corresponding (homeomorphic) fibre of $X \to \pi_0(X)$. Since a point in a w-local space is closed if and only if it is closed in its connected component (which is closed), we have $x \in X^c$, so $y \in Y^c$. Thus $Y^c$ is exactly the set of closed points of $Y$. Moreover, since $Y^c \cong T$, each connected component contains a unique closed point, and for each $y \in Y^c$ the set of generalizations of $y$ is precisely the fibre of $Y \to \pi_0(Y)$ (as any generalization of $y$ lies in the same connected component with $y$). The lemma follows.
\end{proof}

\section{w-local rings}

\begin{definition}[w-local rings, {\cite[Definition 2.2.1(i)]{proetale}}]
    A ring \(A\) is \emph{w-local} if Spec(A) is w-local.
    \label{def:w-local-ring}
    \uses{def:w-local-space}
    \lean{IsWLocalRing}
    \leanok
\end{definition}

\begin{definition}[w-local ring maps, {\cite[Definition 2.2.1(iii)]{proetale}}]
    A ring map \(f: A \to B\) between w-local rings is \emph{w-local} if the induced map of w-local spaces \(\Spec(f) : \Spec(B) \to \Spec(A)\) is w-local.
    % Spectral is automatic, only sending closed pts to closed pts matters.
    \label{def:w-local-ring-map}
    \uses{def:w-local-ring, def:w-local-space, def:w-local-space-map}
\end{definition}

\begin{proof}
  Omitted.
\end{proof}

\begin{lemma}
  \label{thm:surjective-w-local-ring}
  If \(A \to B\) is a surjective ring map and \(A\) is w-local, then \(B\) is w-local.
  \uses{def:w-local-ring}
  \lean{IsWLocalRing.of_surjective}
  \leanok
\end{lemma}

\begin{proof}
  \leanok
  \uses{thm:closed-subspace-w-local}
  This is \Cref{thm:closed-subspace-w-local}.
\end{proof}

\begin{lemma}
  \label{thm:w-local-set-decomposition}
  \uses{def:w-local-ring}
  Let $A$ be a w-local ring. Then $\Spec (A)$ is set theoretically the disjoint union of the spectra of the local
  rings at closed points.
\end{lemma}

\begin{proof}
  This follows from the fact that every point specializes to a unique closed point and the spectrum of the local ring at a point is the set of all points specialize to it.
\end{proof}

\begin{lemma}
  \label{thm:isom-of-identifies-local-rings-bijective}
  \uses{def:identify-local-rings}
  \lean{RingHom.BijectiveOnStalks.bijective_of_bijective}
  \leanok
  Let $A \to B$ be a ring map such that
  \begin{enumerate}
    \item $A \to B$ identifies local rings, and
    \item $\Spec (B) \to \Spec (A)$ is bijective.
  \end{enumerate}
  Then $A \to B$ is an isomorphism.
\end{lemma}

\begin{proof}
    \leanok
    \uses{thm:localization-isom-of-going-down}
    Since $A \to B$ is flat (can be checked on stalks of \(B\)), it has going down. % Algebra.HasGoingDown.of_flat
    Thus \Cref{thm:localization-isom-of-going-down} shows for any prime $\q \subset B$ lying over $\p \subset A$ we have $B_{\q} = B_{\p B}$. Since $B_{\q} = A_{\p}$ by assumption, we see that $A_{\p} = B_{\p B}$ for all primes $\p$ of $A$. Thus $A = B$ since isomorphism can be checked locally on $A$.
\end{proof}

\begin{lemma}[\stacksproject{097E}]
  \label{thm:isom-of-identifies-local-rings-w-local-isom-pi0}
  \lean{RingHom.IsWLocal.bijective_of_bijective}
  \leanok
  \uses{def:identify-local-rings,def:w-local-ring,def:w-local-ring-map}
  Let $A \to B$ be a w-local ring map between w-local rings such that
  \begin{enumerate}
    \item $A \to B$ identifies local rings, and
    \item $\pi_0(\Spec (B)) \to \pi_0(\Spec (A))$ is bijective.
  \end{enumerate}
  Then $A \to B$ is an isomorphism.
\end{lemma}

\begin{proof}
  \leanok
  \uses{thm:closed-points-isom-pi0,thm:w-local-set-decomposition,
    thm:isom-of-identifies-local-rings-bijective}
  By \Cref{thm:isom-of-identifies-local-rings-bijective}, it suffices to show
  that $Y = \Spec (B) \to X = \Spec (A)$ is bijective.
  Let $X^c \subset X$ and $Y^c \subset Y$ be the sets of closed points. By assumption $Y^c$ maps into $X^c$ and the induced map $Y^c \to X^c$ is a bijection since \Cref{thm:closed-points-isom-pi0} holds. By \Cref{thm:w-local-set-decomposition}, we see that both $X$ and $Y$ , as sets, are disjoint unions of the spectra of the local rings at closed points. As \(A \to B\) identifies local rings, $Y \to X$ is a bijection.
\end{proof}

The goal of the remaining part of this section is to show that for every ring $A$, there exists a faithfully flat, ind-Zariski
$A$-algebra $A_{w}$ with $A_w$ w-local. Let $A$ be a ring.

\begin{definition}
    Let $f \in A$ be an element and $I \subseteq A$ an ideal.
    Let $S_{f, I} \subseteq A$ be the multiplicative subset of elements which map to invertible
    elements of $(A/I)_{f}$.
    We define $\tildering{A}{I}{f}$ to be the $A$-algebra $S_{f, I}^{-1}$
    and $\tildeideal{I}{f} \subseteq \tildering{A}{I}{f}$ to be the kernel of
    the induced map $\tildering{A}{I}{f} \to (A / I)_{f}$.
    \label{def:tilde-locclosed}
    \lean{WLocalization.Generalization.submonoid, WLocalization.Generalization, WLocalization.Generalization.ideal}
\end{definition}

If $Z \subseteq \spec{A}$ is a locally closed subscheme of the form $D(f) \cap V(I)$, we
also write $A_{\widetilde{Z}}$ instead of $\tildering{A}{I}{f}$.
In the following, when we use the notation $A_{\widetilde{Z}}$ there is always a canonical
choice of $I$ and $f$ for $Z$.
In general this notation is justified by the fact that, up to isomorphism, the ring $A_{\widetilde{Z}}$ does not depend
on the choice of $f$ and $I$ \stacksproject{096V}.

\begin{lemma}[{\stacksproject{096V}}]
    \uses{def:generalization, def:tilde-locclosed}
    \label{lemma:locally-closed-specializations}
    Let $f \in A$ be an element and $I \subseteq A$ an ideal.
    \begin{enumerate}
        \item The map $\spec{\tildering{A}{I}{f}} \to \spec{A}$ induces
            a homeomorphism onto the generalization of $D(f) \cap V(g)$.
        \item The closed subscheme $V(\tildeideal{I}{f})$ induces an isomorphism
            onto $D(f) \cap V(I)$.
            \label{item:closed-subscheme-tilde}
        \item If $A'$ is a ring with an element $f' \in A$ and ideal $I' \subseteq A'$ and
            $A \to A'$ is a ring map that maps
            $D(f') \cap V(I')$ into $D(f) \cap V(I)$, then there is a unique $A$-algebra map
            $\tildering{A}{I}{f} \to \tildering{A'}{I'}{f'}$ making the obvious
            diagram with $A \to A'$ commute.
    \end{enumerate}
    In particular, every
    point in $\spec{\tildering{A}{I}{f}}$ specializes to a point in $V(\tildeideal{I}{f})$.
    \lean{WLocalization.Generalization.range_algebraMap_generalization,
      WLocalization.Generalization.bijOn_algebraMap_generalization_specComap_zeroLocus_ideal,
      WLocalization.Generalization.submonoid_le,WLocalization.Generalization.map,
      WLocalization.Generalization.exists_specializes_zeroLocus_ideal}
\end{lemma}

\begin{proof}
    \leanok
    Let $S = S_{f, I}$, $Z = D(f) \cap V(I)$ and $J = \tildeideal{I}{f}$.

    Note that $Z$ is homeomorphic to $\spec{(A/I)_f}$.
    Since the map $\spec{S^{-1}A} \to \spec{A}$ is a homeomorphism onto its image, to show (1)
    it suffices to show that the image is the set of points of $\spec{A}$ specializing to $D(f) \cap V(I)$.
    Let $\mathfrak{p}$ be a prime of $A$ specializing to a $\mathfrak{q}$ in $Z$,
    so $I \subseteq \mathfrak{q}$ and $f \not\in \mathfrak{q}$. In particular, the image
    of every element of $\mathfrak{p}$ in $(A / I)_f$ is contained in $\mathfrak{q}$ and
    is therefore not invertible.
    Conversely, let $\mathfrak{p}$ be a prime of $A$ which does not specialize to a point in
    $Z$, hence
    the image of $\mathfrak{p}$ in $(A/I)_f$ contains the unit element. So there exists
    $g \in \mathfrak{p}$ that becomes invertible in $(A / I)_f$, so $g \in S$ and
    $\mathfrak{p}$ is not in the image of $\spec{S^{-1} A}$.

    By construction, the induced map $S^{-1}A \to (A / I)_f$ is surjective and hence
    induces an isomorphism $S^{-1}A / J \to (A / I)_f$. This shows (2).

    Let $g \in A$. Suppose the image of $g$ in $(A' / I')_{f'}$ is not invertible.
    Then there exists a prime ideal $\mathfrak{p}$ of $(A' / I')_{f'}$ that contains the image of $g$.
    Hence $g$ is in the preimage of $\mathfrak{p}$ under $A \to A'$, which lies in $Z$ by assumption. Thus
    $g$ is not invertible in $(A / I)_f$, which shows $g \not\in S$.
\end{proof}

In view of \ref{lemma:locally-closed-specializations}, given $Z = D(f) \cap V(I)$ we may
view $Z$ as a closed subscheme of $\spec{A_{\widetilde{Z}}}$ defined by the ideal
$V( \tildeideal{I}{f})$.

\begin{lemma}
    Let $f \in A$ be an element and $I \subseteq A$ an ideal. Then
    $\tildering{A}{I}{f}$ is ind-Zariski over $A$.
    \uses{def:tilde-locclosed}
    \label{lemma:tilde-locclosed-ind-zariski}
    \lean{WLocalization.Generalization.indZariski}
\end{lemma}

\begin{proof}
    \uses{lemma:ind-Zariski-localization}
    $\tildering{A}{I}{f} = S^{-1} A$ for some multiplicative subset $A$.
\end{proof}

\begin{definition}[Stratification by subsets]
    Let $E, F \subseteq A$ be subsets. We define
    \[
    Z(E, F) = \bigcap_{f \in E'} D(f) \cap \bigcap_{f \in F} V(f)
    .\]
    \label{def:subset-stratum}
    \lean{WLocalization.stratum}
\end{definition}

\begin{lemma}
    Let $A$ be a ring, $E, F \subseteq A$ finite subsets. Then
    \[
    Z(E, F) = D\left( \prod_{f \in E}  f \right) \cap V\left( \sum_{f \in F} fA \right)
    .\]
    \uses{def:subset-stratum}
    \label{lemma:subset-stratum-equals-inter}
    \lean{WLocalization.stratum_eq_basicOpen_inter_zeroLocus}
\end{lemma}

\begin{proof}
    This is immediate from the definitions of $D$ and $V$.
\end{proof}

\begin{lemma}
    Let $E_1, F_1, E_2, F_2 \subseteq A$ be subsets such that $E_1 \subseteq E_2$ and
    $F_1 \subseteq F_2$. Then $Z(E_2, F_2) \subseteq Z(E_1, F_1)$.
    \label{lemma:subset-stratum-mono}
    \lean{WLocalization.stratum_anti}
\end{lemma}

\begin{proof}
    This is immediate from the definition of $Z(-, -)$.
\end{proof}

\begin{lemma}
    Let $A$ be a ring, $E \subseteq A$ a finite subset. Then
    \[
        \spec{A} = \coprod_{E = E' \coprod E''} Z(E', E'')
    \] set theoretically.
    \uses{def:subset-stratum}
    \label{lemma:subset-stratification}
    \lean{WLocalization.Stratification.Index.iUnion_stratum}
\end{lemma}

\begin{proof}
    \leanok
    Let $x \in \spec{A}$ and $E = E' \coprod E''$ a disjoint union decomposition. Then
    $x \in Z(E', E'')$ if and only if $f(x) \neq 0$ for all $f \in E'$ and $f(x) = 0$
    for all $f \in E''$. Hence $x$ is contained in
    \[
    Z(\{ f \in E  \mid f(x) \neq 0\}, \{f \in E  \mid f(x) = 0\})
    ,\] and this is the only set of this form in which $x$ is contained.
\end{proof}

\begin{definition}
    Let $A$ be a ring and $E \subseteq A$ a finite subset. Let
    \[
    A_{E} = \prod_{E = E' \coprod E''} A_{\widetilde{Z(E', E'')}}
    .\]
    Further, we set $I_E \subseteq A$ to be the product of the ideals
    defining $Z(E', E'')$ in $\spec{A_{\widetilde{Z(E', E'')}}}$.
    \uses{def:subset-stratum, def:tilde-locclosed, lemma:subset-stratum-equals-inter}
    \label{def:finite-stratification-ring}
    \lean{WLocalization.ProdStrata, WLocalization.ProdStrata.ideal}
\end{definition}

Note that \ref{def:finite-stratification-ring} is well-defined, because for finite
subsets $E', E''$ of $A$, the locally closed subset $Z(E', E'')$ is always
of the right form by \ref{lemma:subset-stratum-equals-inter}.

\begin{lemma}
    Let $E \subseteq A$ be a finite subset. Then
    $V(I_E) = \coprod_{E = E' \coprod E''} Z(E', E'')$. In particular,
    the map $\spec{A_E} \to \spec{A}$ induces a bijection $V(I_E) \to \spec{A}$.
    \label{lemma:finite-stratification-ideal-bijection}
    \uses{def:finite-stratification-ring}
    \lean{WLocalization.ProdStrata.bijOn_algebraMap_specComap_zeroLocus_ideal}
\end{lemma}

\begin{proof}
    \leanok
    \uses{lemma:subset-stratification}
    By definition $I_E$ is the product of the ideals defining $Z(E', E'')$ in
    $\spec{A_{\widetilde{Z(E', E'')}}}$. Hence the first claim. The second
    follows from the first, because $\spec{A} = \coprod_{E = E' \coprod E''} Z(E', E'')$ by
    \ref{lemma:subset-stratification}.
\end{proof}

\begin{lemma}
    Let $E \subseteq A$ be a finite subset. Then every point of $\spec{A_E}$ specializes
    to a point in $V(I_E)$. In particular, $V(I_E)$ contains all closed points of $\spec{A_E}$.%
    \uses{def:finite-stratification-ring}
    \label{lemma:finite-stratification-closed-points}
    \lean{WLocalization.ProdStrata.exists_specializes_zeroLocus_ideal,
      WLocalization.ProdStrata.mem_zeroLocus_ideal_of_isClosed}
\end{lemma}

\begin{proof}
    \leanok
    \uses{lemma:locally-closed-specializations}
    Let $x \in \spec{A_E}$. Then $x$ is in $\spec{A_{\widetilde{Z(E', E'')}}}$ for some
    $E = E' \coprod E''$. By
    \ref{lemma:locally-closed-specializations}, $x$ specializes to a point
    in $Z(E', E'') \subseteq V(I_E)$.
    If $x$ is closed, then $x$ specializes to itself, so $x \in V(I_E)$.
\end{proof}

Let $A$ be a ring. Given finite subsets $E_1, E_2 \subseteq A$ with $E_1 \subseteq E_2$, there is a
canonical transition map $A_{E_1} \to A_{E_2}$ of $A$-algebras: For
a disjoint union decomposition $E_2 = E_2' \coprod E_2''$, we set
$E_1' = E_1 \cap E_2'$ and $E_1'' = E_1 \cap E_2''$. By \ref{lemma:subset-stratum-mono},
we obtain $Z(E_2', E_2'') \subseteq Z(E_1', E_1'')$ and hence
by \ref{lemma:locally-closed-specializations} a unique $A$-algebra map
$A_{\widetilde{Z(E_1', E_1'')}} \to A_{\widetilde{Z(E_2', E_2'')}}$.

\begin{lemma}
    Let $E_1, E_2 \subseteq A$ be finite subsets with $E_1 \subseteq E_2$.
    The induced map $\spec{A_{E_2}} \to \spec{A_{E_1}}$ maps
    $V(I_{E_2})$ into $V(I_{E_1})$.
    \uses{def:finite-stratification-ring, lemma:subset-stratum-mono}
    \label{lemma:finite-stratification-map-closed}
    \lean{WLocalization.ProdStrata.mapsTo_map_specComap}
\end{lemma}

\begin{proof}
    \leanok
    This holds because for every decomposition $E_2 = E_2' \coprod E_2''$
    and $E_1' = E_1 \cap E_2'$ and $E_1'' = E_1 \cap E_2''$, the induced map
    $\spec{A_{\widetilde{Z(E_2', E_2'')}}} \to \spec{A_{\widetilde{Z(E_1', E_1'')}}}$
    maps $Z(E_2', E_2'')$ into $Z(E_1', E_1'')$.
\end{proof}

\begin{definition}
    Let $I(A)$ be the partially ordered set of all finite subsets of $A$. We define the
    $A$-algebra $A_{w}$ as
    \[
    A_w = \colim_{E \in I(A)} A_E
    .\]
    We denote by $I_w \subseteq A_w$ the colimit of the ideals $I_E$, i.e.,
    the union of the images of the $I_E$ in $A_w$.%
    \uses{def:finite-stratification-ring, lemma:subset-stratum-mono, lemma:finite-stratification-map-closed}
    \label{def:w-localization}
    \lean{WLocalization.commAlgCat,WLocalization,WLocalization.ideal}
\end{definition}

By the following lemma, we see $V(I_w)$ is the inverse limit of the closed subsets $V(I_E)$.
\begin{lemma}
  \label{thm:colim-ideal-lim-zero-locus}
    Let $(A_i, I_i)_{i \in I}$ be a filtered system of
    rings with ideals. Set $A = \colim_{i \in I} A_i$ and
    $I = \colim_{i \in I} I_i = \colim_{i \in I} I_i A$. Then
    \[
    V(I) = \lim_{i \in I} V(I_i) = \bigcap_{i \in I} \text{preimage of } V(I_i) \text{ in } \spec{A}
    .\]
  \lean{zeroLocus_eq_iInter_specComap_preimage,specComap_preimage_zeroLocus_subset}
\end{lemma}

\begin{proof}
  Omitted.
\end{proof}

\begin{lemma}
    The $A$-algebra $A_{w}$ is ind-Zariski.
    \uses{def:w-localization}
    \label{lemma:w-localization-ind-Zariski}
    \lean{WLocalization.indZariski}
\end{lemma}

\begin{proof}
    \uses{lemma:ind-Zariski-products, lemma:ind-ind-Zariski}
    For $E \subseteq A$ a finite subset, the $A$-algebra $A_E$ is ind-Zariski
    as a finite product of ind-Zariski algebras by \ref{lemma:ind-Zariski-products}.
    Thus $A_w$ is a filtered colimit of ind-Zariski algebras, hence ind-Zariski
    by \ref{lemma:ind-ind-Zariski}.
\end{proof}

\begin{lemma}
    \uses{def:w-localization}
    \label{lemma:w-localization-closed-points}
    \lean{WLocalization.bijOn_algebraMap_specComap_zeroLocus_ideal,
      WLocalization.exists_mem_zeroLocus_ideal_specializes,WLocalization.zeroLocus_ideal_eq}
    Let $A$ be a ring.
    \begin{enumerate}
        \item The composition $V(I_w) \to \spec{A_w} \to \spec{A}$ is a bijection.
            \label{item:w-localization-bijection}
        \item Every point of $\spec{A_w}$ specializes to a unique point of $V(I_w)$. \label{item:w-localization-specializes}
        \item $V(I_w)$ is the set of closed points of $\spec{A_w}$. \label{item:w-localization-closed-points}
    \end{enumerate}
\end{lemma}

\begin{proof}
    \uses{thm:colim-ideal-lim-zero-locus,lemma:finite-stratification-ideal-bijection,lemma:finite-stratification-closed-points,lemma:closure-limit-intersection,thm:closure-limit-compatible}
    Since $V(I_w) = \lim_{E} V(I_E)$ (by \Cref{thm:colim-ideal-lim-zero-locus}) and $V(I_E) \to \spec{A_E} \to \spec{A}$ is
    bijective for all $E$ by \Cref{lemma:finite-stratification-ideal-bijection}, the
    first claim holds.

    We first show that every point of $\spec{A_w}$ specializes to a point in $V(I_w)$.
    For this let $y \in \spec{A_w}$. For all $E \subseteq A$ finite, denote by
    $T_E$ the closure of the image of $y$ in $\spec{A_E}$. By
    \Cref{lemma:finite-stratification-closed-points}, the intersection $T_E \cap V(I_E)$ is
    non-empty. Since both $T_E$ and $V(I_E)$ are closed, the intersection is closed
    and hence quasi-compact. Thus $(\lim_E T_E) \cap V(I_w) = \lim_E (T_E \cap V(I_E))$ is non-empty.
    Apply \Cref{lemma:closure-limit-intersection,thm:closure-limit-compatible} to singleton \({y}\), the subset $\lim_E T_E$ is the closure of $y$ in $\spec{A_w}$. (Shuold be a cofiltered version of \verb`IsCompact.nonempty_iInter_of_directed_nonempty_isCompact_isClosed`.) Hence
    $y$ specializes to a point in $V(I_E)$.

    For uniqueness, suppose $y \in \spec{A_w}$ specializes to $z, z' \in V(I_w)$ with
    $z \neq z'$. Denote the images of $z, z'$ in $\spec{A}$ by $x, x'$.
    Because $V(I_w) \to \spec{A_w} \to \spec{A}$ is injective, $x$ and $x'$ are distinct. Hence we may
    assume there exists $f \in A$ such that $x \in D(f)$ and $x' \in V(f)$. Set
    $E = \{f\}$. Then
    \[
        \spec{A_E} = \spec{\tildering{A}{0}{f}} \coprod \spec{\tildering{A}{(f)}{1}}
    .\]
    Denote the images of $y, z, z'$ in $\spec{A_E}$ by $y_E, x_E, x'_E$. Since
    $x_E$ maps to $x \in D(f)$ and $x_E'$ maps to $x' \in V(f)$, they lie
    in different components of the disjoint union decomposition above. But
    since $\spec{A_w} \to \spec{A}$ is continuous, $y_E$ specializes to both $x_E$ and $x'_E$
    which is not possible.

    Finally, every closed point specializes to itself and is hence contained in $V(I_{w})$. Conversely,
    every point in $V(I_{w})$ specializes only to itself, so it is closed.
\end{proof}

\begin{corollary}
    Let $A$ be a ring. The $A$-algebra $A_w$ is w-local.
    \uses{def:w-localization,def:w-local-ring}
    \label{cor:w-localization-w-local}
    \lean{WLocalization.isWLocalRing}
\end{corollary}

\begin{proof}
  \uses{lemma:w-localization-closed-points}
  Immediate from \Cref{item:w-localization-specializes,item:w-localization-closed-points} in \Cref{lemma:w-localization-closed-points}.
\end{proof}

\begin{lemma}
    The $A$-algebra $A_w$ is faithfully flat.
    \uses{def:w-localization}
    \label{lemma:w-localization-faithfully-flat}
    \lean{WLocalization.faithfullyFlat}
\end{lemma}

\begin{proof}
    \leanok
    \uses{lemma:w-localization-closed-points, lemma:w-localization-ind-Zariski}
    By \ref{lemma:w-localization-ind-Zariski}, $A \to A_w$ is flat. Because
    $V(I_w)$ surjects onto $\spec{A}$ by \ref{lemma:w-localization-closed-points},
    in particular $\spec{A_w} \to \spec{A}$ is surjective, hence
    $A \to A_w$ is faithfully flat.
\end{proof}

\begin{lemma}[\stacksproject{00GA}]
    Let $B$ be an $A$-algebra and $\mathfrak{m}$ a maximal ideal of of $R$. Let
    $\mathfrak{q}$ be a prime ideal lying over $\mathfrak{m}$ such that
    $\kappa(\mathfrak{q})$ is algebraic over $\kappa(\mathfrak{m})$. Then
    $\mathfrak{q}$ is maximal.
    \label{lemma:prime-over-maximal-algebraic}
    \lean{Ideal.IsMaximal.of_isAlgebraic}
    \leanok
\end{lemma}

\begin{proof}
    \leanok
    We have a factorization $\kappa(\mathfrak{m}) = R / \mathfrak{m} \to S / \mathfrak{q} \to \kappa(\mathfrak{q})$.
    Hence $S / \mathfrak{q}$ is an intermediate subring of an algebraic field extension and therefore
    a field.
\end{proof}

\begin{lemma}
    Let $A$ be w-local and $I \subseteq A$ an ideal. Then
    $A_{\widetilde{V(I)}} = A_{\widetilde{(I,1)}}$ is w-local. If $V(I)$ only contains closed points, then
    the set of closed points is $V(IA_{\widetilde{V(I)}})$.
    \label{lemma:w-local-tilde-w-local}
    \lean{isWLocalRing_generalization_one,IsWLocalRing.zeroLocus_map_algebraMap_generalization_one_eq}
\end{lemma}

\begin{proof}
    \uses{lemma:generalization-specialization-closed, lemma:locally-closed-specializations,
      lemma:w-local-specialization-closed}
    By \ref{lemma:locally-closed-specializations}, $\spec{A_{\widetilde{V(I)}}}$ is homeomorphic
    to the generalization $V(I)^g$. Since $\spec{A}$ is w-local and $V(I)$ is closed
    under specialization, also $V(I)^g$ is closed under specialization by
    \ref{lemma:generalization-specialization-closed}. Hence
    $\spec{A_{\widetilde{V(I)}}} \cong V(I)^g$ is w-local by \ref{lemma:w-local-specialization-closed}.

    By \Cref{item:closed-subscheme-tilde} in \Cref{lemma:locally-closed-specializations}, we may identify $V(IA_{\widetilde{V(I)}})$ with
    $V(I)$.
    Since every point in
    $\spec{A_{\widetilde{V(I)}}} \cong V(I)^g$ specializes to a unique point in $V(I)$,
    the closed subset $V(I)$ contains all closed points of $\spec{A_{\widetilde{V(I)}}}$. Hence,
    if $V(I)$ only contains closed points, the closed points are exactly $V(I)$.
\end{proof}

\begin{proposition}[Universal property of w-localization, \stacksproject{0977}]
    Let $A \to B$ be a ring map with $B$ w-local. There exists a unique
    factorization $A \to A_w \to B$ such that $A_w \to B$ is w-local.
    \uses{def:w-localization, def:w-local-ring-map}
    \label{lemma:univ-prop-w-localization}
\end{proposition}

\begin{proof}
    TBA.
\end{proof}

\begin{definition}[w-localization of ring map]
    Let $f\colon A \to B$ be a ring map. We denote by $f_w\colon A_w \to B_w$ the
    unique w-local ring map making
    \[
    \begin{tikzcd}
        A \arrow{r} \arrow{d} & B \arrow{d} \\
        A_w \arrow{r} & B_w
    \end{tikzcd}
    \] commute.
    \label{def:w-localization-map}
    \uses{lemma:univ-prop-w-localization}
\end{definition}

\begin{lemma}
    Let $A \to B$ be a ring map that induces algebraic extensions on residue fields
    and let $I \subseteq A$ be an ideal, such that $V(I)$ only contains closed points.
    Then $V(IB)$ only contains closed points.
    \label{lemma:preimage-closed-points-algebraic}
    \lean{zeroLocus_map_algebraMap_subset_closedPoints}
    \leanok
\end{lemma}

\begin{proof}
    \uses{lemma:prime-over-maximal-algebraic}
    \leanok
    By \ref{lemma:prime-over-maximal-algebraic}, every prime of $B$ lying above a maximal ideal of $A$
    is maximal. In particular, the closed subset $V(I B) = (\spec{B} \to \spec{A})^{-1}(V(I))$ only contains
    closed points.
\end{proof}

\begin{lemma}
    Let $B$ be an $A$-algebra satisfying going down. If every closed point of $\spec{A}$
    has a preimage in $\spec{B}$, the map $\spec{B} \to \spec{A}$ is surjective.
    \label{lemma:going-down-surjective-of-closed}
    \lean{Algebra.HasGoingDown.specComap_surjective_of_closedPoints_subset_preimage}
    \leanok
\end{lemma}

\begin{proof}
    \leanok
    Let $\mathfrak{p}$ be a prime of $A$. Then it is contained in a maximal ideal $\mathfrak{m}$
    that is the preimage of a prime $\mathfrak{q}$ of $B$. By going-down, there
    exists therefore a prime $\mathfrak{q}' \subseteq \mathfrak{q}$ lying above $\mathfrak{p}$.
\end{proof}

\begin{definition}
  Let $A$ be a ring and $I \subseteq A$ an ideal. We define the \emph{w-localization of $A$ with respect to $I$} to be the ring $(A_w)_{\widetilde{V(IA_{w})}}$, denoted by $A_{w,I}$.
  \uses{def:w-localization,def:tilde-locclosed}
  \label{def:w-localization-ideal}
  \lean{Ideal.WLocalization}
  \leanok
\end{definition}

\begin{lemma}
  Let $A$ be a ring and $I \subseteq A$ an ideal. The ring $A_{w,I}$ is w-local.
  \uses{def:w-localization-ideal,def:w-local-ring}
  \label{lemma:w-localization-ideal-w-local}
  \lean{isWLocalRing_generalization_one,Ideal.WLocalization.isWLocalRing}
\end{lemma}

\begin{proof}
  \uses{lemma:w-local-tilde-w-local, cor:w-localization-w-local}
  Follows immediately from \Cref{lemma:w-local-tilde-w-local} and \Cref{cor:w-localization-w-local}.
\end{proof}

\begin{lemma}
  Let $A$ be a ring and $I \subseteq A$ an ideal. The ring map $A \to A_{w, I}$ is ind-Zariski.
  \label{thm:w-localization-ideal-ind-Zariski}
  \uses{def:w-localization-ideal,def:ind-Zariski}
  \lean{Ideal.WLocalization.indZariski}
  \leanok
\end{lemma}

\begin{proof}
  \leanok
  \uses{lemma:w-localization-ind-Zariski, lemma:tilde-locclosed-ind-zariski,thm:ind-Zariski-composition}
  The map $A \to A_w$ is ind-Zariski by \Cref{lemma:w-localization-ind-Zariski} and
  the map $A_w \to A_{w, I}$ is ind-Zariski by \Cref{lemma:tilde-locclosed-ind-zariski}.
  Hence the composition $A \to A_{w, I}$ is ind-Zariski by \Cref{thm:ind-Zariski-composition}.
\end{proof}

\begin{lemma}
    \uses{def:w-localization-ideal}
    \label{lemma:closed-closed-points-tilde-w-local}
    \lean{Ideal.WLocalization.zeroLocus_map_algebraMap_eq_closedPoints,
      Ideal.WLocalization.quotientMap_algebraMap_bijective,
      Ideal.WLocalization.bijOn_zeroLocus_map}
    Let $I \subseteq A$ be an ideal such that $V(I) \subseteq A$ only contains closed points. Then
    \begin{enumerate}
        \item the set of closed points of $A_{w, I}$ is $V(IA_{w, I})$, and
        \item the map $A / I \to A_{w, I} / I A_{w, I}$ is an isomorphism. \label{item:closed-closed-points-tilde-w-local-quot-isom}
    \end{enumerate}
\end{lemma}

\begin{proof}
    \uses{lemma:w-localization-ind-Zariski,thm:ind-Zariski-identifies-local-rings,lemma:w-local-tilde-w-local,thm:isom-of-identifies-local-rings-bijective,lemma:preimage-closed-points-algebraic,lemma:tilde-locclosed-ind-zariski,lemma:locally-closed-specializations}
    $A \to A_w$ is ind-Zariski by \ref{lemma:w-localization-ind-Zariski} and therefore induces
    algebraic extensions of residue fields by \ref{thm:ind-Zariski-identifies-local-rings}.
    In particular, the closed subset $V(I A_w)$ only contains
    closed points by \ref{lemma:preimage-closed-points-algebraic} and
    the first claim follow from \ref{lemma:w-local-tilde-w-local}.

    The map $A \to A_{w, I}$ is ind-Zariski as a composition of ind-Zariski maps by
    \ref{lemma:w-localization-ind-Zariski} and \ref{lemma:tilde-locclosed-ind-zariski}.
    Since ind-Zariski is stable under base change, the same holds for $A/I \to A_{w, I}/IA_{w, I}$.
    Hence it identifies local rings. By \Cref{item:closed-subscheme-tilde} in \Cref{lemma:locally-closed-specializations},
    $V(IA_{w, I}) \to V(I A_w)$ is an isomorphism. Moreover
    $V(I A_{w}) \to V(I)$ is a bijection by \Cref{item:w-localization-bijection}
    because $V(I A_{w}) \subseteq V(I_w)$. Hence, the composition
    $\spec{A_{w, I} / I A_{w, I}} = V(I A_{w, I}) \to V(I A_w) \to V(I)$ is bijective. Thus
    $A / I \to A_{w, I} / I A_{w, I}$ is an isomorphism by \ref{thm:isom-of-identifies-local-rings-bijective}.
\end{proof}

\begin{lemma}[{\cite[\href{https://stacks.math.columbia.edu/tag/097A}{Tag 097A}]{stacks-project}}, (2)]
\label{thm:closed-points-isom-w-local}
\uses{def:w-local-ring,def:w-local-ring-map,def:ind-Zariski,def:w-localization-ideal}
\lean{Ideal.WLocalization.algebraMap_specComap_preimage_closedPoints_eq}
\leanok
Let $A$ be a w-local ring. Let $I \subset A$ be an ideal cutting out the set $X^c$ of closed points in $X = \Spec (A)$.
Let $A \to B$ be a ring map inducing algebraic extensions on residue fields at primes. Let $C$ be the ring $B_{w, IB}$ constructed in \Cref{def:w-localization-ideal}. Then
\begin{enumerate}
    \item $B/IB \to C/IC$ is an isomorphism, \label{item:quotinet-isom-closed-points-isom-w-local} % thus surjective on closed points
    \item $C$ is w-local, \label{item:w-local-closed-points-isom-w-local}
    \item the map $p: Y = \Spec (C) \to X$ induced by the ring map $A \to B \to C$ is w-local, and % this is not needed?
    \item $p^{-1}(X^c)$ is the set of closed points of $Y$, i.e. V(IC) is the set of closed points of $Y$. \label{item:inverse-image-closed-points-isom-w-local}
\end{enumerate}
\end{lemma}

\begin{proof}
    \leanok
    \uses{lemma:w-localization-closed-points,lemma:closed-closed-points-tilde-w-local,
      lemma:tilde-locclosed-ind-zariski,def:w-localization-ideal,lemma:w-localization-ideal-w-local}
    Every point of $V(IB)$ is a closed point of $\spec{B}$ by 
    \Cref{lemma:preimage-closed-points-algebraic}.
    Hence we may choose $C = (B_w)_{\widetilde{V(IB_w)}}$ as in 
    \Cref{def:w-localization-ideal}. \Cref{item:quotinet-isom-closed-points-isom-w-local} holds
    by \Cref{lemma:closed-closed-points-tilde-w-local}. 
    \Cref{item:w-local-closed-points-isom-w-local} holds by 
    \Cref{lemma:w-localization-ideal-w-local}. Then
    the set of closed points of $C$ is $V(IC)$, which is the preimage of $V(I)$ under
    $\spec{C} \to \spec{B} \to \spec{A}$.
\end{proof}

\begin{lemma}
  Let $A$ be a w-local ring. Let $I \subset A$ be an ideal cutting out the set $X^c$ of closed points in $X = \Spec (A)$.
  Let $B$ be a faithfully flat $A$-algebra. Let $C$ be the ring $B_{w, IB}$ constructed in \Cref{def:w-localization-ideal}. Then
  the composition $A \to B \to C$ is also faithfully flat.
  \label{cor:closed-points-isom-w-local-faithfully-flat}
  \uses{def:w-localization-ideal}
  \lean{Ideal.WLocalization.faithfullyFlat_map_algebraMap}
  \leanok
\end{lemma}

\begin{proof}
  \leanok
  \uses{thm:ind-Zariski-is-flat,thm:w-localization-ideal-ind-Zariski,lemma:going-down-surjective-of-closed,lemma:closed-closed-points-tilde-w-local}
  By \Cref{thm:ind-Zariski-is-flat,thm:w-localization-ideal-ind-Zariski}, $A \to B \to C$ is flat, so by \ref{lemma:going-down-surjective-of-closed} it suffices
  to show every closed point of $A$ is in the image. But since $\spec{B} \to \spec{A}$ is surjective,
  also $V(IC) \cong V(IB) \to V(I)$ is surjective (the first map is an isomorphism by \Cref{item:closed-closed-points-tilde-w-local-quot-isom} in \Cref{lemma:closed-closed-points-tilde-w-local}).
\end{proof}


\section{w-strictly-local Rings}

\begin{definition}[w-strictly local rings, {\cite[Definition 2.2.1(ii)]{proetale}}]
    A ring \(A\) is \emph{w-strictly local} if \(A\) is w-local, and every local ring of \(A\) at a
    maximal ideal is strictly henselian.
    \label{def:w-strictly-local-ring}
    \uses{def:w-local-ring,def:strictly-henselian-local-ring}
    \lean{IsWStrictlyLocalRing}
    \leanok
\end{definition}

The goal of this section is to show that every ring $A$ admits a faithfully flat, ind-étale algebra $B$ that is $w$-strictly-local.

\begin{lemma}
  \label{thm:ind-etale-strictly-henselian-localization-isom}
  \uses{def:ind-etale,def:strictly-henselian-local-ring,def:identify-local-rings}
  Let \(A\) be a strictly henselian local ring. Let \(A \to B\) be an ind-étale ring map.
  Let $\n$ be an maximal ideal of $B$ lying over the maximal ideal $\m$ of $A$.
  Then the map $A \to B_\n$ is isomorphism.
\end{lemma}

\begin{proof}
  \uses{thm:etale-over-strictly-henselian-localization-isom}
  Write $B = \colim_i B_i$ as a filtered colimit of étale $A$-algebras $B_i$. Since localization commutes with colimits, we have $B_\n = \colim_i (B_i)_{\n_i}$ where $\n_i$ is the restriction of $\n$ to $B_i$. Then use \Cref{thm:etale-over-strictly-henselian-localization-isom} to conclude each $B_i$ is isomorphic to $A$.
\end{proof}

\begin{lemma}
    Let $A$ be a ring such that every faithfully flat étale ring map $A \to B$ has
    a retraction. Then every local ring of $A$ at a maximal ideal is strictly henselian.
    \label{lemma:retractions-strictly-henselian}
\end{lemma}

\begin{proof}
    Let $\mathfrak{m}$ be a maximal ideal of $A$, denote by $\kappa$ the residue field $\kappa(\mathfrak{m})$
    and by $\kappa^{\mathrm{sep}}$ a separable closure of $\kappa$. Let
    $A_{\mathfrak{m}} \to B \to \kappa^{\mathrm{sep}}$ be a factorisation of $A_{\mathfrak{m}} \to \kappa^{\mathrm{sep}}$
    where $A_{\mathfrak{m}} \to B$ is étale. Since $A_{\mathfrak{m}} = \colim_{f \not\in \mathfrak{m}} A_f$,
    by \ref{lemma:etale-ind-spreads} there exists $f \not\in \mathfrak{m}$ and an étale $A_f$-algebra $B'$
    such that $B = A_{\mathfrak{m}} \otimes_{A_f} B'$.
    %Denote by $\mathfrak{q}$ the kernel of
    %the composition $B' \to B \to \kappa^{\mathrm{sep}}$.
    %Then $\mathfrak{q}$ lies over $\mathfrak{m}$ in $A$.
    Because $A \to A_f \to B'$ is étale,
    there are finitely many prime ideals $\mathfrak{q}_1, \ldots, \mathfrak{q}_n$ of $B'$ lying over $\mathfrak{m}$.
    Since the kernel of the composition $B' \to B \to \kappa^{\mathrm{sep}}$ lies over $\mathfrak{m}$
    (the kernel of $A \to \kappa_{\mathrm{sep}}$ is $\mathfrak{m}$), we have $n \ge 1$.
    Set $\mathfrak{q} = \mathfrak{q}_1$. By prime avoidance, there exists $g \in B'$ such that
    $g \in \mathfrak{q}_i$ for all $i \ge 2$ and $g \not\in \mathfrak{q}$. Hence
    $\mathfrak{q}$ is the unique prime ideal of $B'_g$ lying over $\mathfrak{m}$.

    Let $U$ be the image of the induced map $\spec{B'_g} \to \spec{A}$. Since $A \to B'$ is étale, $U$ is open.
    Since $\{\mathfrak{m}\}$ is closed in $\spec{A}$, there exist $a_i$ such that
    \[
        \spec{A} - \{\mathfrak{m}\} = \bigcup_{i} D(a_i)
    .\] The complement of $U$ is closed and quasicompact, and contained in $\spec{A} - \{\mathfrak{m}\}$, hence
    there exist $a_1, \ldots, a_r$ such that
    \[
        \spec{A} = U \cup \bigcup_{i=1}^r D(a_i)
    .\] Hence the map $A \to B'_g \times \prod_{i=1}^{r} A_{a_i}$ is étale and faithfully flat. By assumption,
    it therefore has a retraction $\sigma$. Since $\mathfrak{q}$ is the unique prime ideal of $B'_g \times \prod_{i=1}^{r} A_{a_i}$
    lying over $\mathfrak{m}$, we have $\sigma^{-1}(\mathfrak{m}) = \mathfrak{q}$. In particular,
    we obtain a retraction
    \[
        B'_{\mathfrak{q}} = \left(B'_g \times \prod_{i=1}^{r} A_{a_i}\right)_{\mathfrak{q}} \to A_{\mathfrak{m}}
    \] of the map $A_{\mathfrak{m}} \to B'_{\mathfrak{q}}$. Precomposing with $B' \to B'_{\mathfrak{q}}$ we obtain
    a map $B' \to A_{\mathfrak{m}}$ and hence a map
    $B = B' \otimes_{A_f} A_{\mathfrak{m}} \to A_{\mathfrak{m}}$ that is a retraction of $A_{\mathfrak{m}} \to B$.
\end{proof}

\begin{definition}
    Let $A$ be a ring. Denote by $S(A)$ the set of finite subsets of faithfully flat étale $A$-algebras.
    We call the $A$-algebra $T(A) = \colim_{E \in S(A)} \bigotimes_{B \in E} B$ the \emph{étale precontraction} of
    $A$.
    \label{def:etale-precontraction}
    \lean{CategoryTheory.MorphismProperty.precontraction}
\end{definition}

\begin{lemma}
    \uses{def:etale-precontraction}
    \label{lemma:etale-precontraction-semi-retraction}
    For every faithfully flat étale ring map $A \to B$, there exists an $A$-algebra map $B \to T(A)$.
    \lean{CategoryTheory.MorphismProperty.Precontraction.π}
\end{lemma}

\begin{proof}
    $\{B\}$ is an element of $S(A)$, hence $B$ itself is part of the diagram defining $T(A)$.
\end{proof}

\begin{lemma}
    Let $R$ be a commutative ring.
    \begin{itemize}
        \item The forgetful functors from commutative $R$-algebras to $R$-algebras and
            from $R$-algebras to $R$-modules preserve filtered colimits.
        \item The forgetful functor from commutative $R$-algebras to commutative rings
            preserves filtered colimits.
        \item Tensoring on the left with an $R$-algebra preserves colimits on the
            category of commutative $R$-algebras.
    \end{itemize}
    \label{lemma:commalgcat-colimits}
    \lean{CommAlgCat.preservesColimitsOfShape_tensorLeft,
        CommAlgCat.preservesFilteredColimitsOfSize_forget_commRingCat,
        CommAlgCat.preservesFilteredColimitsOfSize_forget_algCat,
        AlgCat.preservesFilteredColimitsOfSize_forget_moduleCat}
\end{lemma}

\begin{lemma}
    Let $B = \colim_i B_i$ be a filtered colimit of (faithfully) flat $A$-algebras. Then
    $B$ is (faithfully) flat over $A$.
    \label{lemma:ind-faithfully-flat}
    \lean{Module.Flat.of_colimitPresentation,Module.FaithfullyFlat.of_colimitPresentation,
      Module.Flat.iff_ind_flat,Module.FaithfullyFlat.iff_ind_faithfullyFlat,
      RingHom.Flat.iff_ind_flat,RingHom.FaithfullyFlat.iff_ind_faithfullyFlat,
      RingHom.Flat.of_isColimit,RingHom.FaithfullyFlat.of_isColimit}
\end{lemma}

\begin{proof}
    \uses{lemma:commalgcat-colimits}
    Because colimits commute with tensor products, ind-flat is flat. One also checks
    that pro-surjective is surjective, hence the claim.
\end{proof}

\begin{lemma}
    $T(A)$ is ind-étale and faithfully flat over $A$.
    \label{lemma:etale-precontraction-ind-etale}
    \uses{def:etale-precontraction, def:ind-etale}
    \lean{CategoryTheory.MorphismProperty.pro_precontraction_hom}
\end{lemma}

\begin{proof}
    \uses{lemma:ind-faithfully-flat}
    For every $E \in S(A)$, the $A$-algebra $\bigotimes_{B \in E} B$ is étale and faithfully flat. Hence
    the result follows from the definition of ind-étale and \Cref{lemma:ind-faithfully-flat}.
\end{proof}

\begin{definition}
    Let $A$ be a ring. Set $T^0(A) = A$ and for $n \in \N$ set $T^{n+1}(A) = T(T^n(A))$.
    We call the $A$-algebra $T^{\infty}(A) = \colim_{n \in \N} T^n(A)$
    the \emph{étale contraction of $A$}.
    \uses{def:etale-precontraction}
    \label{def:etale-contraction}
    \lean{IndEtaleContraction}
    \leanok
\end{definition}

\begin{lemma}
    $T^{\infty}(A)$ is ind-étale and faithfully flat over $A$.
    \uses{def:etale-contraction,def:ind-etale}
    \label{lemma:etale-contraction-ind-etale}
    \lean{faithfullyFlat_indEtaleContraction, indEtale_indEtaleContraction}
\end{lemma}

\begin{proof}
    \uses{lemma:etale-precontraction-ind-etale, lemma:ind-faithfully-flat, lemma:ind-ind-etale}
    This follows from \Cref{lemma:etale-precontraction-ind-etale}, \Cref{lemma:ind-faithfully-flat}
    and \Cref{lemma:ind-ind-etale}.
\end{proof}

\begin{proposition}[{\stacksproject{097R}}]
    Let $T^{\infty}(A) \to B$ be a faithfully flat and étale $A$-algebra map. Then
    it has a retraction.
    \uses{def:etale-contraction}
    \label{prop:etale-contraction-retraction}
    \lean{RingHom.Etale.exists_comp_eq_id_indContraction}
    \leanok
\end{proposition}

\begin{proof}
    \leanok
    \uses{lemma:etale-ind-spreads, lemma:etale-precontraction-semi-retraction}
    By \Cref{lemma:etale-ind-spreads}
    there exists $n \in \N$ and a faithfully flat, étale $T^n(A)$-algebra $B'$ such that
    $B = T^{\infty}(A) \otimes_{T^n(A)} B'$. By \Cref{lemma:etale-precontraction-semi-retraction},
    there exists a map $B' \to T^{n+1}(A)$. The identity of $T^{\infty}(A)$ and the composition
    $B' \to T^{n+1}(A) \to T^{\infty}(A)$ now induce a retraction
    \[
    B = T^{\infty}(A) \otimes_{T^n(A)} B' \to T^{\infty}(A)
    .\]
\end{proof}

\begin{corollary}
    Let $\mathfrak{m}$ be a maximal ideal of $T^{\infty}(A)$. Then $T^{\infty}(A)_{\mathfrak{m}}$
    is strictly Henselian.
    \uses{def:etale-contraction}
    \label{cor:strictly-henselian-etale-contraction}
\end{corollary}

\begin{proof}
    \uses{prop:etale-contraction-retraction, lemma:retractions-strictly-henselian}
    This follows from \Cref{prop:etale-contraction-retraction} and \Cref{lemma:retractions-strictly-henselian}.
\end{proof}

\begin{lemma}
    Let $k$ be a field and $A$ an ind-étale, local $k$-algebra. Then $A$ is
    a separable algebraic field extension of $k$.
    \uses{def:ind-etale}
    \label{lemma:field-local-ind-etale}
    \lean{Algebra.IndEtale.isSeparable}
\end{lemma}

\begin{proof}
    \leanok
    Let $A = \colim_i A_i$ such that $A_i$ is an étale $k$-algebra for all $i$.
    Let $\mathfrak{m}$ be the maximal ideal of $A$. Since localization commutes with colimits,
    we obtain $A = A_{\mathfrak{m}} = \colim_i (A_i)_{\mathfrak{m}_i}$ where
    $\mathfrak{m}_i$ is the restriction of $\mathfrak{m}$ to $A_i$. Since $A_i$ is étale over $k$,
    it is a finite product of finite separable extensions of $k$, hence $(A_i)_{\mathfrak{m}_i}$
    is a finite separable extension of $k$. Since a filtered colimit of finite separable
    extensions is an algebraic separable extension, we conclude.
\end{proof}

\begin{proposition}
    Let $f\colon A \to B$ be ind-étale. Then $f$ induces separable algebraic 
    extensions on residue fields.
    \uses{def:ind-etale}
    \label{prop:ind-etale-residue-fields}
    \lean{Algebra.IndEtale.isSeparable_residueField}
\end{proposition}

\begin{proof}
    \uses{lemma:field-local-ind-etale}
    \leanok
    Let $\mathfrak{q}$ be a prime ideal of $B$ lying over $\mathfrak{p} \subseteq A$.
    Write $B = \colim_i B_i$ as a filtered colimit of étale $A$-algebras. Every element
    $s \in B$ lifts to some $s_i \in B_i$, and the composition
    \[
        \kappa(\mathfrak{p}) \otimes_A B_i \to \kappa(\mathfrak{q})
    \]
    is an algebra homomorphism from the étale $\kappa(\mathfrak{p})$-algebra
    $\kappa(\mathfrak{p}) \otimes_A B_i$ to the local ring $\kappa(\mathfrak{q})$, so by
    \ref{lemma:field-local-ind-etale} the image of $s$ in $\kappa(\mathfrak{q})$ is separable
    over $\kappa(\mathfrak{p})$. A general element of $\kappa(\mathfrak{q})$ is a ratio of two
    such images, and separability passes to ratios via the separable closure.
\end{proof}

\begin{proposition}[{\cite[\href{https://stacks.math.columbia.edu/tag/097X}{Tag 097X}]{stacks-project}}]
    For any ring \(A\), there exists an ind-étale faithfully flat \(A\)-algebra \(B\) with \(B\) w-strictly local.
    \label{thm:ind-etale-w-strictly-local-cover}
    \uses{def:w-strictly-local-ring,def:ind-etale}
\end{proposition}

\begin{proof}
    \uses{cor:strictly-henselian-etale-contraction, prop:ind-etale-residue-fields,def:w-localization-ideal,thm:closed-points-isom-w-local,cor:closed-points-isom-w-local-faithfully-flat,
      lemma:w-localization-closed-points, lemma:etale-contraction-ind-etale}
    By \Cref{lemma:w-localization-closed-points}, the ring $A_w$ is w-local, so the set of closed points
    of $\spec{A_w}$ is closed.
    Let $I$ be the radical ideal of $A_w$ such that $V(I)$ is the set of closed points of $\spec{A_w}$.
    Set $B' = T^{\infty}(A_w)$.
    By \Cref{cor:strictly-henselian-etale-contraction}, the local ring $B'_{\mathfrak{m}}$ is
    strictly henselian for every maximal ideal $\mathfrak{m}$ of $B'$.
    The map $A_w \to B'$ induces algebraic extensions on residue fields at primes by
    \Cref{lemma:etale-contraction-ind-etale} and \Cref{prop:ind-etale-residue-fields},
    so applying \Cref{def:w-localization-ideal,thm:closed-points-isom-w-local} to $A_w \to B'$, we obtain
    an ind-Zariski ring map $B' \to B$ such that $B'/IB' \to B/IB$ is an isomorphism,
    $V(IB)$ is the set of closed points of $\spec{B}$, $B$ is w-local
    and $A_w \to B' \to B$ is faithfully flat by \Cref{cor:closed-points-isom-w-local-faithfully-flat}. Since ind-Zariski maps identify
    local rings by \Cref{thm:ind-Zariski-identifies-local-rings} and every closed point of
    $\spec{B}$ maps to a closed point of $B'$,
    the local rings at maximal ideals of $B$ are strictly Henselian. Hence $B$ is w-strictly local.

    It remains to verify that the composition $A \to A_w \to B' \to B$ is ind-étale and faithfully flat.
    This follows because both ind-étale and faithfully flat is stable under composition.
\end{proof}

\begin{remark}
    By following all constructions, we may choose $B$
    as $((T^{\infty}(A_w))_w)_{\widetilde{V(I ((T^{\infty}(A_w))_w))}}$
    where $I$ is the radical ideal of $A_w$ defining the closed points of $A_w$.

    One may ask why $T^\infty(A)_w$ does not work. The reason for this is
    that the map $\spec{T^{\infty}(A)_w} \to \spec{T^{\infty}(A)}$ does not necessarily
    map closed points to closed points. If $A$ is already $w$-local,
    we can eliminate the points mapping to non-closed points by taking $I$ to be the ideal defining
    the closed points of $\spec{A}$.
    Then $V(I T^{\infty}(A))$ only contains closed points and by
    passing to the generalization of the preimage of $V(I T^{\infty}(A))$ in
    $\spec{T^{\infty}(A)_w}$, we obtain the desired affine scheme. If $A$ is not $w$-local,
    we may replace $A$ by $A_w$ in this process and obtain the result.
\end{remark}

\begin{definition}[{\cite[\href{https://stacks.math.columbia.edu/tag/097C}{Tag 097C}]{stacks-project}}]
  \label{def:colim-open-closed-localizations}
  \lean{WContractification.RestrictClopen.map,WContractification.Restriction.diag,
    WContractification.Restriction}
  Let $A$ be a ring, $X = \Spec(A)$, and $T \subset \pi_0(X)$. Define the ring $A_T$ by
  \[ 
    A_T := \colim_{W} A_W,
  \]
  where $W$ runs over all open and closed subsets of $X$ containing the inverse image $Z$ of $T$, 
  and $A \to A_W$ is the localization of $A$ at $W$.
\end{definition}

\begin{lemma}[{\cite[\href{https://stacks.math.columbia.edu/tag/097C}{Tag 097C}]{stacks-project}}]
  \label{thm:colim-open-closed-localizations-surj-ind-zariski}
  \uses{def:colim-open-closed-localizations,def:ind-Zariski}
  \lean{WContractification.Restriction.indZariski,
    WContractification.Restriction.algebraMap_surjective,
    WContractification.Restriction.range_algebraMap_specComap,
    WContractification.Restriction.isClosedEmbedding_algebraMap_specComap}
  Let $A$ be a ring, $X = \Spec(A)$, and $T \subset \pi_0(X)$ a closed subset. 
  Then the map $A \to A_T$ from \Cref{def:colim-open-closed-localizations} is a surjective 
  ind-Zariski ring map, and the induced morphism $\Spec(A_T) \to \Spec(A)$ restricts to a 
  homeomorphism of $\Spec(A_T)$ onto the inverse image of $T$ in $X$.
\end{lemma}

\begin{proof}
  \leanok
  \uses{thm:intersection-closed-open-iff}
  Each $A \to A_W$ is a localization at an idempotent. Passing to filtered colimits, $A \to A_T$ is surjective and ind-Zariski.

  When $T$ is closed, its inverse image $Z$ is a closed union of connected components. By \Cref{thm:intersection-closed-open-iff}, we have
  \[
    Z = \bigcap_{\substack{Z \subseteq W \\ W \text{ open and closed}}} W.
  \]
  Thus, $\Spec(A_T) = \varprojlim_W W$ is homeomorphic to $Z$ via the natural map.
\end{proof}

In the following series of constructions, we will modify the connected components of a spectrum to match a given profinite space as in {\cite[\href{https://stacks.math.columbia.edu/tag/097D}{Tag 097D}]{stacks-project}}.

\begin{definition}
  \label{def:modify-pi0-finite}
  \uses{def:colim-open-closed-localizations,thm:pi0-prod,thm:pi0-isom}
  \lean{WContractification.Pullback}
  \leanok
  Let $A$ be a ring and let $X = \Spec (A)$. Let $T$ be profinite space and let 
  $f : T \to \pi_0(X)$ be a continuous map. Let \(S\) be a finite discrete quotient of 
  \(T\) with quotient map \(q : T \to S\).  Define the ring $A^f_{q}$ as follow: 
  let $Z = \operatorname{Im} (T \to \pi_0(X) \times S = \pi_0(\Spec (A^{S})))$, 
  where \(A^{S} = \prod_{s \in S} A\) and let 
  \[A^f_{q} := (A^{S})_{Z}\] 
  be the $A^S$-algebra as in \Cref{def:colim-open-closed-localizations}.
\end{definition}

\begin{lemma}
  \label{thm:modify-pi0-finite-properties}
  \uses{def:modify-pi0-finite,def:ind-Zariski,def:identify-local-rings}
  \lean{WContractification.Pullback.indZariski,
        WContractification.Pullback.bijectiveOnStalks_algebraMap}
  Let $A$ be a ring and let $X = \Spec (A)$. Let $T$ be profinite space and let $f : T \to \pi_0(X)$ be a continuous map. Let \(S\) be a finite discrete quotient of \(T\) with quotient map \(q : T \to S\). Then the $A^S$-algebra defined in \Cref{def:modify-pi0-finite} viewed as $A$-algebra through $A \to A^S \to A^f_{q}$ is ind-Zariski. In particular, it identifies local rings. The map $A^S \to A^f_{q}$ induces a homeomorphism of $\Spec (A^f_{q})$ onto \((X \times S) \times_{\pi_0(X) \times S} Z = X \times_{\pi_0(X)} Z\).
\end{lemma}

\begin{proof}
  \uses{thm:ind-Zariski-identifies-local-rings,thm:colim-open-closed-localizations-surj-ind-zariski,thm:ind-Zariski-composition,thm:ind-Zariski-identifies-local-rings}
  The map $A \to A^S$ is ind-Zariski by \Cref{lemma:ind-Zariski-products} and $A^S \to A^f_{q}$ is also ind-Zariski by \Cref{thm:colim-open-closed-localizations-surj-ind-zariski}.  The composition $A \to A^f_{q}$ is therefore ind-Zariski as well by \Cref{thm:ind-Zariski-composition}. By \Cref{thm:ind-Zariski-identifies-local-rings}, $A \to A^f_{q}$ identifies local rings. The last sentence is a direct consequence of \Cref{thm:colim-open-closed-localizations-surj-ind-zariski}.
\end{proof}

\begin{definition}
  \label{def:modify-pi0-transition-map}
  \lean{WContractification.PullbackProfinite.diag}
  \uses{def:modify-pi0-finite,def:identify-local-rings,def:identifies-local-ring-hom-set-bijection}
  \leanok
  Let $A$ be a ring and let $X = \Spec (A)$. Let $T$ be profinite space and let $f : T \to \pi_0(X)$ be a continuous map. Let \(S\) and \(S'\) be finite discrete spaces and let \(q : T \to S\) and \(q' : T \to S'\) be quotient maps. Let \(g : S' \to S\) be a map such that \(q = g \circ q'\). Let \(Z \subseteq \pi_0(\Spec (A^{S})) \times S\) and \(Z' \subseteq \pi_0(\Spec (A^{S'})) \times S'\) be the images of \(T\) under the maps induced by \(f\) and \(f'\) respectively. Then \(g\) induces a map \(Z \to Z'\) and a map \(\tilde{g} : \Spec (A^f_{q}) \cong X \times_{\pi_0(X)} Z \to X \times_{\pi_0(X)} Z' \cong \Spec (A^{f'}_{q'})\).

  Since both maps \(A \to A^f_{q}\) and \(A \to A^{f'}_{q'}\) identify local rings, by the bijection in \Cref{def:identifies-local-ring-hom-set-bijection} we can find a transition ring map
  \[t_{g}: A^f_{q} \to A^{f'}_{q'}\]
  that induces the topological map \(\tilde{g}\) between the spectra.
\end{definition}

\begin{lemma}
  \label{thm:modify-pi0-transition-map-compatible}
  \lean{WContractification.PullbackProfinite.diag}
  \uses{def:modify-pi0-transition-map}
  \leanok
  Let $A$ be a ring and let $X = \Spec (A)$. Let $T$ be profinite space and let $f : T \to \pi_0(X)$ be a continuous map. Let \(S\), \(S'\) and \(S''\) be finite discrete spaces and let \(q : T \to S\), \(q' : T \to S'\) and \(q'' : T \to S''\) be quotient maps. Let \(g : S' \to S\) and \(g' : S'' \to S'\) be maps such that \(q = g \circ q'\) and \(q' = g' \circ q''\). Then the transition maps defined in \Cref{def:modify-pi0-transition-map} satisfy
  \[t_{g'} \circ t_{g} = t_{g \circ g'}.\]
\end{lemma}

\begin{proof}
  \uses{thm:modify-pi0-finite-properties,def:identifies-local-ring-hom-set-bijection}
  \leanok
  Both maps \(t_{g'}\) and \(t_{g} \circ t_{g' \circ g}\) induce the same topological map between the spectra by construction. Since all maps involved identify local rings by \Cref{thm:modify-pi0-finite-properties}, by fully faithfulness in \Cref{def:identifies-local-ring-hom-set-bijection} we conclude that they are equal.
\end{proof}

\begin{definition}\leanok
[{\cite[\href{https://stacks.math.columbia.edu/tag/097D}{Tag 097D}]{stacks-project}}]
  \label{def:modify-pi0-profinite}
  \lean{WContractification.PullbackProfinite}
  \uses{def:modify-pi0-finite,def:modify-pi0-transition-map,thm:modify-pi0-transition-map-compatible}
  Let $A$ be a ring and let $X = \Spec (A)$. Let $T$ be a profinite space and let $T \to \pi_0(X)$ be a continuous map. Write $T = \lim T_i$ as the limit of an inverse system of finite discrete spaces over a directed set. Define the ring \(A^f_{\pi_0}\) as follows:
  \[
    A^f_{\pi_0} := \colim_{(T \to T_i)} A^f_{q_i},
  \]
  where \(q_i : T \to T_i\) are the quotient maps and \(A^f_{q_i}\) are the rings defined in \Cref{def:modify-pi0-finite} and the transition maps are those defined in \Cref{def:modify-pi0-transition-map}. By \Cref{thm:modify-pi0-transition-map-compatible} the transition maps are compatible thus the colimit is well-defined.

  On the Lean side the indexing category is \texttt{(DiscreteQuotient T)\^{}op}: Mathlib's order \(S_1 \leq S_2\) on \texttt{DiscreteQuotient T} encodes refinement (\(S_1\) finer than \(S_2\), with \(\top\) the coarsest one-class quotient), so a Mathlib morphism \(S_1 \to S_2\) is the finer-to-coarser direction at the type level. The blueprint's transition \(t_g : A^f_q \to A^{f'}_{q'}\) going from a coarser \(q\) to a finer \(q'\) at the algebra level therefore corresponds to a covariant functor on \texttt{(DiscreteQuotient T)\^{}op}, and the colimit \(A^f_{\pi_0}\) is the colimit of that functor in \texttt{CommAlgCat A}.
\end{definition}

\begin{lemma}[{\cite[\href{https://stacks.math.columbia.edu/tag/097D}{Tag 097D}]{stacks-project}}]
  \label{thm:modify-pi0-profinite-properties}
  \lean{WContractification.PullbackProfinite.indZariski, WContractification.PullbackProfinite.faithfullyFlat, WContractification.PullbackProfinite.isWLocalRing, WContractification.PullbackProfinite.pi0Equiv}
  \uses{def:modify-pi0-profinite,def:ind-Zariski}
  Let $A$ be a ring and let $X = \Spec (A)$. Let $T$ be a profinite space and let $T \to \pi_0(X)$ be a continuous map. Then the $A$-algebra $A^f_{\pi_0}$ defined in \Cref{def:modify-pi0-profinite} is ind-Zariski. Let $Y = \Spec (A^f_{\pi_0})$. Then $\pi_0(Y) \cong T$ and the diagram
  \[
  \begin{tikzcd}
  Y \arrow[r] \arrow[d] & T \arrow[d] \\
  X \arrow[r] & \pi_0(X)
  \end{tikzcd}
  \]
  is cartesian in the category of topological spaces.
\end{lemma}

\begin{proof}
  \uses{thm:modify-pi0-finite-properties,lemma:ind-ind-Zariski,thm:cartesian-profinite}
  By \Cref{thm:modify-pi0-finite-properties} we see that each \(A \to A^f_{q_i}\) is ind-Zariski, thus $A^f_{\pi_0}$ is ind-Zariski by \Cref{lemma:ind-ind-Zariski}. We have the following canonical cartesian diagram
  \[
  \begin{tikzcd}
  \Spec (A^f_{q_i}) \arrow[r] \arrow[d] & Z_i \arrow[d] \\
  X \arrow[r] & \pi_0(X).
  \end{tikzcd}
  \]
  Because $T = \lim Z_i$ and, $Y = \Spec (A^f_{\pi_0}) = \lim \Spec (A^f_{q_i})$ fits into the cartesian diagram
  \[
  \begin{tikzcd}
  Y \arrow[r] \arrow[d] & T \arrow[d] \\
  X \arrow[r] & \pi_0(X)
  \end{tikzcd}
  \]
  of topological spaces. By \Cref{thm:cartesian-profinite} we conclude that $T = \pi_0(Y)$.
\end{proof}

\begin{lemma}[{\cite[\href{https://stacks.math.columbia.edu/tag/09AZ}{Tag 09AZ}, second part in the proof of (3) $\implies$ (1)]{stacks-project}}]
  \label{thm:ff-identifies-local-rings-plus-c-has-retraction-if}
  \uses{def:w-local-ring,def:extremally-disconnected,def:identify-local-rings}
  Let $A$ be a ring. If the following two conditions hold:
  \begin{enumerate}
    \item $A$ is w-local, and
    \item $\pi_0(\operatorname{Spec}(A))$ is extremally disconnected.
  \end{enumerate}
  Then every faithfully flat ring map $A \to B$ identifying local rings with \(B\) w-local and whose closed points of \(\Spec (B)\) are exactly \(V(IB)\) has a retraction.
\end{lemma}

\begin{proof}
  \uses{thm:closed-points-isom-pi0,thm:extremally-disconnected-projective,def:colim-open-closed-localizations,thm:colim-open-closed-localizations-surj-ind-zariski,thm:surjective-w-local-ring,thm:ind-Zariski-identifies-local-rings,thm:w-local-map-composition,thm:identifies-local-rings-composition,thm:isom-of-identifies-local-rings-w-local-isom-pi0}
  Let $A \to B$ be faithfully flat and identifying local rings, with \(B\) w-local and whose closed points of \(\Spec (B)\) are exactly \(V(IB)\). We will show that $A \to B$ has a retraction.

  % Let $I \subset A$ be the radical ideal such that $V(I) \subset \Spec(A)$ is the set of closed points of $\Spec (A)$. We may replace $B$ by the ring $C$ constructed in \Cref{thm:closed-points-isom-w-local} for $A \to B$ and $I \subset A$. The map \(B \to C\) is faithfully flat and ind-Zariski, thus identifies local rings by \Cref{thm:ind-Zariski-identifies-local-rings}. Thus we may assume $\Spec (B)$ is w-local such that the set of closed points of $\Spec (B)$ is $V(IB)$ by \Cref{item:w-local-closed-points-isom-w-local,item:inverse-image-closed-points-isom-w-local} in \Cref{thm:closed-points-isom-w-local}.

  Choose a continuous section to the surjective continuous map $V(IB) \to V(I)$. This is possible because $V(I) = \Spec(A)^c \cong \pi_0(\Spec(A))$ (by \Cref{thm:closed-points-isom-pi0}) is extremally disconnected, see \Cref{thm:extremally-disconnected-projective}. The image is a closed subspace $T \subset \pi_0(\Spec(B)) \cong V(IB)$ (being the continuous image of the quasi-compact space $V(I)$) that maps homeomorphically onto $\pi_0(\Spec(A))$. 
  
  Let $B \to B_T$ be the ring map from \Cref{def:colim-open-closed-localizations}, which is surjective and ind-Zariski by \Cref{thm:colim-open-closed-localizations-surj-ind-zariski}. It also induces a homeomorphism $\pi_0(\Spec(B_T)) \cong T \cong \pi_0(\Spec(A))$.
  Moreover, by \Cref{thm:surjective-w-local-ring}, $B_T$ is w-local and $B \to B_T$ is w-local. Since ind-Zariski maps identify local rings (\Cref{thm:ind-Zariski-identifies-local-rings}), the composition $A \to B \to B_T$ remains w-local (\Cref{thm:w-local-map-composition}) and identifies local rings (\Cref{thm:identifies-local-rings-composition}). Therefore, by \Cref{thm:isom-of-identifies-local-rings-w-local-isom-pi0}, $A \to B_T$ is an isomorphism.
\end{proof}

\section{w-contractible rings}

\begin{definition}[w-contractible rings, {\cite[Definition 2.4.1]{proetale}}]
    A ring \(A\) is \emph{w-contractible} if it is w-strictly local and $\pi_0(\spec{A})$
    is extremally disconnected.
    \label{def:w-contractible-ring}
    \lean{IsWContractibleRing}
    \leanok
\end{definition}

We will see that if a ring $A$ is w-contractible, every faithfully flat, ind-étale map $A \to B$ has a retraction. To show
this we will perform a series of reductions.

\begin{lemma}
    \label{thm:ind-etale-plus-c-has-retraction-if-w-contractible}
    \uses{def:w-contractible-ring,def:extremally-disconnected,def:ind-etale}
    \lean{IsWContractibleRing.exists_retraction_of_zeroLocus_map_eq_closedPoints}
    \leanok
    Let $A$ be w-contractible and $I \subseteq A$ an ideal cutting out the set $X^c$ of closed points in $X = \spec{A}$.
    Then every faithfully flat ind-étale map \(A \to B\) with \(B\) w-local and
    whose closed points of \(\Spec (B)\) are exactly \(V(IB)\) has a retraction.
\end{lemma}

\begin{proof}
  \uses{thm:ind-etale-base-change,thm:ind-etale-strictly-henselian-localization-isom,thm:ff-identifies-local-rings-plus-c-has-retraction-if}
  \leanok
  Let \(A \to B\) be a faithfully flat ind-étale map with \(B\) w-local and such that the closed points of \(\Spec(B)\) are exactly \(V(IB)\). In this case, \(A \to B\) identifies local rings. To justify this, it suffices to check the condition at the maximal ideals of \(B\), which lie over maximal ideals of \(A\) since they map to \(V(I) = (\Spec(A))^c\). Since \(A\) is w-strictly local, the local rings of \(A\) at its maximal ideals are strictly Henselian.

  Now fix a maximal ideal \(\mathfrak{n}\) of \(B\) and let \(\mathfrak{m} = \mathfrak{n} \cap A\) be its preimage in \(A\), which is a maximal ideal. By base change to \(A_{\mathfrak{m}}\) and applying \Cref{thm:ind-etale-base-change}, we reduce to the case where \(B\) is ind-étale over a strictly Henselian local ring \(A\), with a maximal ideal \(\mathfrak{n}\) lying over the maximal ideal \(\mathfrak{m}\) of \(A\). We then want to show \(B_{\mathfrak{n}} \cong A\). This follows from \Cref{thm:ind-etale-strictly-henselian-localization-isom}.

  Having established that \(A \to B\) identifies local rings, we conclude by \Cref{thm:ff-identifies-local-rings-plus-c-has-retraction-if} that it admits a retraction.
\end{proof}

\begin{proposition}[{\cite[\href{https://stacks.math.columbia.edu/tag/0982}{Tag 0982}, (3) $\implies$ (2)]{stacks-project}}] % actually iff, but we only need this direction
  \label{thm:w-contractible-retraction}
  \uses{def:w-contractible-ring,def:w-strictly-local-ring,def:extremally-disconnected}
  If $A$ is w-contractible, every ind-étale, faithfully flat map $A \to B$ has a retraction.
  \lean{IsWContractibleRing.exists_retraction}
  \leanok
\end{proposition}

\begin{proof}
  \uses{def:w-localization-ideal,cor:closed-points-isom-w-local-faithfully-flat,
    thm:w-localization-ideal-ind-Zariski,thm:ind-Zariski-is-ind-etale,
    thm:ind-etale-comp,thm:closed-points-isom-w-local,
    thm:ind-etale-plus-c-has-retraction-if-w-contractible}
  \leanok
  With \Cref{thm:ind-etale-plus-c-has-retraction-if-w-contractible}, it suffices to reduce to 
  the case that \(B\) is furthur w-local and the closed points of \(\Spec (B)\) 
  are exactly \(V(IB)\).

  Let $A \to B$ be faithfully flat and ind-étale. We may replace $B$ by the ring $C = B_{w, IB}$ 
  (see \Cref{def:w-localization-ideal}). To justify this, note that the map \(A \to C\) is 
  faithfully flat by \Cref{cor:closed-points-isom-w-local-faithfully-flat} and \(B \to C\) is 
  ind-Zariski by \Cref{thm:w-localization-ideal-ind-Zariski}, thus ind-étale by 
  \Cref{thm:ind-Zariski-is-ind-etale}. Therefore, by \Cref{thm:ind-etale-comp}, the composition 
  \( A \to C \) is also ind-étale. Hence, we may assume $\Spec (B)$ is w-local such that the set 
  of closed points of $\Spec (B)$ is $V(IB)$ by \Cref{item:w-local-closed-points-isom-w-local,
  item:inverse-image-closed-points-isom-w-local} in \Cref{thm:closed-points-isom-w-local}.

  % Assume (1) or equivalently (2). We have (3)(c) by Lemma \ref{lemma:097V}. Properties (3)(a) and (3)(b) follow from Lemma \ref{lemma:09AZ}.
\end{proof}

\begin{lemma}
  If a ring \(C\) is w-strictly local, then there exists an ind-Zariski faithfully flat \(C\)-algebra \(D\) with \(D\) w-contractible.
  \label{thm:ind-etale-w-contractible-cover-of-w-strictly-local}
  \uses{def:w-contractible-ring,def:w-strictly-local-ring,def:ind-etale}
  \lean{exists_isWContractibleRing_of_isWStrictlyLocal}
  \leanok
\end{lemma}

\begin{proof}

\leanok
  % The proof goes as in https://stacks.math.columbia.edu/tag/0983, starting from the second sentence.
  \uses{thm:extremally-disconnected-cover,def:modify-pi0-profinite,thm:modify-pi0-profinite-properties,thm:cartesian-profinite-w-local,thm:ind-Zariski-identifies-local-rings}
  Choose an extremally disconnected space $T$ and a surjective continuous map $T \to \pi_0(\operatorname{Spec}(C))$ by \Cref{thm:extremally-disconnected-cover}. (\(T\) can be chosen as \(\beta((\pi_0(\operatorname{Spec}(C)))^{\textnormal{disc}})\).) Note that $T$ is profinite.
  By \Cref{thm:modify-pi0-profinite-properties}, \Cref{def:modify-pi0-profinite} gives an ind-Zariski ring map
  $C \to D$ such that $\pi_0(\Spec (D)) \to \pi_0(\Spec (C))$ realizes $T \to \pi_0(\Spec (C))$ and such that
  \[
  \begin{tikzcd}
  \operatorname{Spec}(D) \arrow[r] \arrow[d] & \pi_0(\operatorname{Spec}(D)) \arrow[d] \\
  \operatorname{Spec}(C) \arrow[r] & \pi_0(\operatorname{Spec}(C))
  \end{tikzcd}
  \]
  is cartesian in the category of topological spaces.

  Following \Cref{thm:cartesian-profinite-w-local}, we note that $\Spec (D)$ is w-local, that $\Spec (D) \to \Spec (C)$ is w-local, and that the set of closed points of $\Spec (D)$ is the inverse image of the set of closed points of $\Spec (C)$.

  Thus it is still true that the local rings of $D$ at its maximal ideals are strictly henselian (as they are isomorphic to the local rings at the corresponding maximal ideals of $C$ by \Cref{thm:ind-Zariski-identifies-local-rings}). Hence $D$ is w-contractible.
\end{proof}

\begin{theorem}[{\cite[Theorem 2.4.9]{proetale}}, {\cite[\href{https://stacks.math.columbia.edu/tag/0983}{Tag 0983}]{stacks-project}}]
  For any ring $A$, there exists an ind-étale faithfully flat $A$-algebra $D$ with $D$ w-contractible.
  \label{thm:ind-etale-w-contractible-cover-exists}
  \uses{def:w-contractible-ring,def:ind-etale}
  \lean{exists_isWContractibleRing}
  \leanok
\end{theorem}

\begin{proof}
    \leanok
    \uses{thm:ind-etale-w-strictly-local-cover,thm:ind-etale-w-contractible-cover-of-w-strictly-local}
    This is a combination of \Cref{thm:ind-etale-w-strictly-local-cover} and \Cref{thm:ind-etale-w-contractible-cover-of-w-strictly-local}.
\end{proof}

\begin{corollary}
    For any ring $A$, there exists an ind-étale faithfully flat $A$-algebra $B$ such that
    every ind-étale faithfully flat map $B \to C$ has a retraction.
    \uses{def:ind-etale}
    \label{cor:ind-etale-retraction-cover-exists}
    \lean{exists_forall_exists_retraction}
    \leanok
\end{corollary}

\begin{proof}
    \leanok
    \uses{thm:ind-etale-w-contractible-cover-exists, thm:w-contractible-retraction}
    This follows from \Cref{thm:ind-etale-w-contractible-cover-exists} and \Cref{thm:w-contractible-retraction}.
\end{proof}
